%\documentclass[12pt,oneside,a4paper]{book}

%\usepackage{amsmath}
%\usepackage{mathtools}
%\usepackage{amsfonts}
%\usepackage{unicode-math}
%\usepackage{geometry}
%\usepackage{ctex}

%\begin{document}

%\setcounter{chapter}{0}
\setcounter{section}{0}
\setcounter{subsection}{0}

\chapter{实数}

\begin{dingyi}
对集合$K$\\
有$|K|\geq 2$且$K$上有二元运算$+,\times:K\times K\to K$满足:\\
对任意$x,y,z\in K$,有$x+(y+z)=(x+y)+z$\\
存在$0\in K$,对任意$x\in K$,有$x+0=0+x=x$\\
对任意$x\in K$,存在$-x\in K$,有$x+(-x)=(-x)+x=0$\\
对任意$x,y\in K$,有$x+y=y+x$\\
对任意$x,y,z\in K$,有$x\times (y\times z)=(x\times y)\times z$\\
存在$1\in K$,对任意$x\in K$,有$x\times 1=1\times x=x$\\
对任意$x\in K,x\neq 0$,存在$\frac{1}{x}\in K$,有$x\times (\frac{1}{x})=(\frac{1}{x})\times x=1$\\
对任意$x,y\in K$,有$x\times y=y\times x$\\
对任意$x,y,z\in K$,有$(x+y)\times z=x\times z+y\times z$\\
则称$K$为数域\\
则称$\Q$为有理数域,则称$\R$为实数域,则称$\C$为复数域
\end{dingyi}

\begin{dingli}
对数域$K$\\
对任意$x\in K$,有$0\times x=x\times 0=0,-(-x)=x$\\
对$x,y,z\in K$,若$x+z=y+z$,则$x=y$\\
对$x,y,z\in K$,若$x\times z=y\times z,z\neq 0$,则$x=y$\\
对$x,y\in K$,有$x\times y=0$当且仅当$x=0$或$y=0$\\
对$x,y\in K$,有$(-x)\times y=-(x\times y),(-x)\times(-y)=x\times y$\\
重要地,有$0\neq 1$
\end{dingli}

\begin{dingyi}
对数域$K$\\
若存在$K_{+}\subset K$和$K$中偏序关系$>$\\
对任意$x\in K$,有$x=0$或$x\in K_{+}$或$-x\in K_{+}$且只有其一成立\\
对任意$x,y\in K_{+}$,有$x+y\in K_{+}$且$x\times y\in K_{+}$\\
对任意$x,y\in K$,有偏序关系$x>y$当且仅当$x+(-y)\in K_{+}$\\
则称$K$为序域\\
则$\Q$为序域,则$\R$为序域,则$\C$不为序域
\end{dingyi}

\begin{zhu}
记$x<y$为$y>x$\\
记$x\geq y$为$x>y$或$x=y$\\
记$x\leq y$为$x<y$或$x=y$
\end{zhu}

\begin{dingli}
对数域$K$为序域\\
则对任意$x,y\in K$,有$x=y$或$x>y$或$y>x$且只有其一成立\\
则对任意$x,y,z\in K$,若$x>y$,则$x+z>y+z$\\
则对任意$x,y,z\in K$,若$x>y,z>0$,则$x\times z>y\times z$
\end{dingli}

\begin{dingli}
对数域$K$为序域\\
则对任意$x\in K$,有$x\times x\geq 0$\\
则对任意$x\in K$,有$x>0$当且仅当$\frac{1}{x}>0$\\
则$1>0$
\end{dingli}

\begin{dingyi}
对数域$K$为序域\\
对$a,b\in K$且$a<b$\\
记开区间为$(a,b)=\{x\in K|a<x<b\}\subset K$\\
记闭区间为$[a,b]=\{x\in K|a\leq x\leq b\}\subset K$\\
记$|a-b|=\sqrt{(a-b)^2}$为$K$上的距离,则$|a-b|=|b-a|$\\
特别地,若$a=b$,则$|b-a|=|a-b|=0$
\end{dingyi}

\begin{dingyi}
对数域$K$为序域\\
对数列$\{x_n\}\subset K$,若存在$x\in K$\\
对任意$\varepsilon>0$,存在$N\in \N$,对任意$n\geq N$,有$|x_n-x|<\varepsilon$\\
则称$\{x_n\}$为收敛的,否则称$\{x_n\}$为发散的\\
若$\{x_n\}$收敛于$x\in K$,则记为$\lim\limits_{n\to\infty}x_n=x$\\
若对任意$M>0$,存在$N\in \N$,对任意$n\geq N$,有$x_n\geq M$,则记为$\lim\limits_{n\to\infty}x_n=\infty$\\
若对任意$M>0$,存在$N\in \N$,对任意$n\geq N$,有$x_n\leq-M$,则记为$\lim\limits_{n\to\infty}x_n=-\infty$
\end{dingyi}

\begin{zhu}
若$\{x_n\}$收敛于$0$,则称$\{x_n\}$为无穷小数列
\end{zhu}

\begin{dingli}
对数域$K$为序域\\
对数列$\{x_n\},\{y_n\}\subset K$\\
若存在$x_0,y_0\in K$,有$\lim\limits_{n\to\infty}x_n=x_0,\lim\limits_{n\to\infty}y_n=y_0$\\
则有$\lim\limits_{n\to\infty}x_n+y_n=x_0+y_0$且$\lim\limits_{n\to\infty}x_n\times y_n=x_0\times y_0$\\
若$y_0\neq 0$且对任意$n\in\N^*$,有$y_n\neq 0$\\
则有$\lim\limits_{n\to\infty}x_n\times\frac{1}{y_n}=x_0\times\frac{1}{y_0}$
\end{dingli}

\begin{dingli}
对数域$K$为序域\\
对数列$\{x_n\},\{y_n\},\{z_n\}\subset K$\\
若对任意$n\in \N^*$,有$x_n\leq z_n\leq y_n$且存在$k\in K$,有$\lim\limits_{n\to\infty}x_n=\lim\limits_{n\to\infty}y_n=k$\\
则有$\lim\limits_{n\to\infty}z_n=x_0=y_0$
\end{dingli}

\begin{dingyi}
对数域$K$为序域\\
若对子集$U\subset K$,对$x\in K$,对任意$u\in U$,有$u\leq x$\\
则称$x$为$U$的上界\\
若对子集$U\subset K$,对$x\in K$,对任意$u\in U$,有$u\geq x$\\
则称$x$为$U$的下界\\
若对子集$U\subset K$,对$U$的上界$\sup U\in K$,对任意$U$的上界$x$,有$\sup U\leq x$\\
则称$\sup U$为$U$的上确界\\
若对子集$U\subset K$,对$U$的下界$\inf U\in K$,对任意$U$的下界$x$,有$\inf U\geq x$\\
则称$\inf U$为$U$的下确界\\
若对子集$U\subset K$,存在$M\in K$,对任意$u\in U$,有$|u|\leq M$\\
则称$U$为有界的
\end{dingyi}

\begin{zhu}
对数列$\{x_n\}\subset K$和给定常数$N\in\N$\\
若存在$M\in K$,对任意$n\in \N^*$,有$|x_n|\leq M$\\
则称$\{x_n\}$为有界的\\
记$\sup\limits_{n\geq N}x_n=\sup\{x_n|n\geq N\}$,特别地,记$\sup x_n=\sup \{x_n\}$\\
记$\inf\limits_{n\geq N}x_n=\inf\{x_n|n\geq N\}$,特别地,记$\inf x_n=\inf \{x_n\}$
\end{zhu}

\begin{dingyi}
对数域$K$为序域\\
对数列$\{x_n\}\subset K$\\
则称$\varlimsup\limits_{n\to\infty}=\lim\limits_{n\to\infty}\sup\limits_{k\geq n}x_k$为$\{x_n\}$的上极限\\
则称$\varliminf\limits_{n\to\infty}=\lim\limits_{n\to\infty}\inf\limits_{k\geq n}x_k$为$\{x_n\}$的下极限
\end{dingyi}

\begin{dingli}
对数域$K$为序域\\
对数列$\{x_n\}\subset K$,有$\varlimsup\limits_{n\to\infty} x_n\geq\varliminf\limits_{n\to\infty}x_n$\\
则有$\varlimsup\limits_{n\to\infty} x_n=\varliminf\limits_{n\to\infty}x_n$当且仅当存在$x_0\in K\cup\{\pm \infty\}$,有$\lim\limits_{n\to\infty}x_n=x_0$
\end{dingli}

\begin{dingyi}
对数域$K$为序域\\
若对任何有上界的子集$U\subset K$,存在$\sup U\in K$,则称$K$为完备的\\
则$\Q$不为完备的,则$\R$为完备的
\end{dingyi}

\begin{dingli}
Archimedes:\\
对自然数集$\N$,则$\N$无上界\\
对任意$\varepsilon>0$,存在$N\in\N$,对任意$n\geq N,n\in \N$,有$\frac{1}{n}<\varepsilon$
\end{dingli}

\begin{dingli}完备Archimedes序域只有$\R$
\end{dingli}

\begin{dingli}
Dedekind:\\
对实数域$\R$\\
对子集$A\neq\varnothing$和$B\neq \varnothing$,若$A\cup B=\R$且对任意$a\in A,b\in B$,有$a\leq b$\\
则存在$c\in\R$,对任意$a\in A,b\in B$,有$a\leq c\leq b$
\end{dingli}

\begin{dingli}
Cantor:\\
对实数域$\R$\\
对数列$\{a_n\}$和$\{b_n\}$,若对任意$n\in \N^*$,有$a_{n}\leq a_{n+1}<b_{n+1}\leq b_{n}$\\
则有闭区间套$[a_1,b_1]\supset[a_2,b_2]\supset\cdots\supset[a_n,b_n]\supset\cdots$\\
若有$\lim\limits_{n\to\infty}(b_n-a_n)=0$,则存在唯一$c\in \R$,有$c\in \bigcap\limits_{n=1}^{\infty}[a_n,b_n]$
\end{dingli}

\begin{dingli}
Heine-Borel:\\
对实数域$\R$\\
对任意闭区间$[a,b]$和开区间族$\{U_{i}|i\in I\}$\\
若$[a,b]\subset \bigcup\limits_{i\in I}U_{i}$,则存在$J\subset I,|J|<\infty$,有$[a,b]\subset \bigcup\limits_{j\in J}U_{j}$
\end{dingli}

\begin{dingli}
Bolzano-Weierstrass:\\
对实数域$\R$\\
对子集$U\subset \R,|U|=\infty$,若$U$为有界的,则存在$x\in\R$为聚点\\
即对任意$\varepsilon>0$,有$|U\cap (x-\varepsilon,x+\varepsilon)|=\infty$\\
对数列$\{x_n\}$,若$\{x_n\}$为有界的\\
则存在子列$\{x_{n_i}\}\subset\{x_n\},x_0\in\R$,有$\lim\limits_{i\to\infty}x_{n_{i}}=x_{0}$
\end{dingli}

\begin{dingli}
Cauchy:\\
对实数域$\R$\\
对数列$\{x_n\}$,若$\{x_n\}$为基本列\\
即对任意$\varepsilon>0$,存在$N\in \N$,对任意$n,m\geq N$,有$|x_n-x_m|<\varepsilon$\\
则存在$x_0\in \R$,有$\lim\limits_{n\to\infty}x_n=x_0$
\end{dingli}

\begin{dingli}
Weierstrass:\\
对实数域$\R$\\
对数列$\{x_n\}$,若$\{x_n\}$为有界的且对任意$n\in \N^*$,有$x_{n+1}\geq x_{n}$\\
则存在$x_0\in \R$,有$\lim\limits_{n\to\infty}x_n=x_0$
\end{dingli}

\begin{dingli}
Bolzano:\\
对实数域$\R$\\
对子集$U\subset\R$,若$U$有上界,则存在$x\in \R$,有$\sup U=x$
\end{dingli}

\begin{dingyi}
对实数域$\R$\\
若有数列$\{a_n\}\subset \R$,记$s_n=\sum\limits_{i=1}^{n}a_i,\sum\limits_{n=1}^{\infty}a_n=\lim\limits_{n\to\infty}s_n$\\
则称$s_n$为$\{a_n\}$的部分和,则称$\sum\limits_{n=1}^{\infty}a_n$为级数\\
若对任意$n\in\N^*$,有$a_n\geq 0$,则称$\sum\limits_{n=1}^{\infty}a_n$为正项级数\\
若$\lim\limits_{n\to\infty}s_n$为收敛的,则称$\sum\limits_{n=1}^{\infty}a_n$为收敛的,记为$|\sum\limits_{n=1}^{\infty}a_n|<\infty$\\
若$\lim\limits_{n\to\infty}s_n$为发散的,则称$\sum\limits_{n=1}^{\infty}a_n$为发散的\\
若$\sum\limits_{n=1}^{\infty}|a_n|$为收敛的,则称$\sum\limits_{n=1}^{\infty}a_n$为绝对收敛的\\
若$\sum\limits_{n=1}^{\infty}a_n$为收敛的但$\sum\limits_{n=1}^{\infty}|a_n|$为发散的,则称$\sum\limits_{n=1}^{\infty}a_n$为条件收敛的
\end{dingyi}

\begin{dingli}
对实数域$\R$\\
有正项级数$\sum\limits_{n=1}^{\infty}a_n$和$\sum\limits_{n=1}^{\infty}b_n$\\
则正项级数$\sum\limits_{n=1}^{\infty}a_n$为收敛的当且仅当存在$M>0$,对任意$n\in\N^*$,有$\sum\limits_{i=1}^{n}a_i\leq M$\\
若$\sum\limits_{n=1}^{\infty}a_n$和$\sum\limits_{n=1}^{\infty}b_n$均为收敛的,则$\sum\limits_{n=1}^{\infty}\sum\limits_{i=1}^{n}a_{i}b_{n+1-i}=\sum\limits_{n=1}^{\infty}a_n\sum\limits_{n=1}^{\infty}b_n$为收敛的
\end{dingli}

\begin{dingli}
对实数域$\R$\\
有正项级数$\sum\limits_{n=1}^{\infty}a_n$和$\sum\limits_{n=1}^{\infty}b_n$\\
d'Alembert:\\
若存在$N\in \N^*$和$r<1$,对任意$n\geq N$,有$\frac{a_{n+1}}{a_{n}}\leq r$,则$\sum\limits_{n=1}^{\infty}a_n$为收敛的\\
若存在$N\in \N^*$和$r\geq 1$对任意$n\geq N$,有$\frac{a_{n+1}}{a_{n}}\geq r$,则$\sum\limits_{n=1}^{\infty}a_n$为发散的\\
Cauchy:\\
若存在$N\in \N^*$和$r<1$,对任意$n\geq N$,有$\sqrt{a_n}\leq r$,则$\sum\limits_{n=1}^{\infty}a_n$为收敛的\\
若存在$r\geq 1$,对无穷多个$n\in\N^*$,有$\sqrt{a_n}\geq r$,则$\sum\limits_{n=1}^{\infty}a_n$为发散的\\
Raabe:\\
若存在$\varepsilon>0,N\in\N^*$,对任意$n\geq N$,有$n\frac{a_n}{a_{n+1}}-(n+1)\geq \varepsilon$,则$\sum\limits_{n=1}^{\infty}a_n$为收敛的\\
若存在$N\in\N^*$,对任意$n\geq N$,有$n\frac{a_n}{a_{n+1}}-(n+1)\leq 0$,则$\sum\limits_{n=1}^{\infty}a_n$为发散的\\
Bertrand:\\
若存在$\varepsilon>0,N\in\N^*$,对任意$n\geq N$,有$n\ln n\frac{a_n}{a_{n+1}}-(n+1)\ln (n+1)\geq \varepsilon$\\
则$\sum\limits_{n=1}^{\infty}a_n$为收敛的\\
若存在$N\in\N^*$,对任意$n\geq N$,有$n\ln n\frac{a_n}{a_{n+1}}-(n+1)\ln (n+1)\leq 0$\\
则$\sum\limits_{n=1}^{\infty}a_n$为发散的\\
Kummmer:\\
若存在正项级数$\sum\limits_{n=1}^{\infty}\frac{1}{c_n}$为发散的,记$K_n=c_n\frac{a_n}{a_{n+1}}-c_{n+1}$\\
若存在$\varepsilon>0,N\in\N^*$,对任意$n\geq N$,有$K_n\geq \varepsilon$.则$\sum\limits_{n=1}^{\infty}a_n$为收敛的\\
若存在$N\in\N^*$,对任意$n\geq N$,有$K_n\leq 0$,则$\sum\limits_{n=1}^{\infty}a_n$为发散的
\end{dingli}

\begin{dingli}
对实数域$\R$\\
若有级数$\sum\limits_{n=1}^{\infty}a_n$和$\sum\limits_{n=1}^{\infty}b_n$均为收敛的且$c,d\in \R$\\
则$\sum\limits_{n=1}^{\infty}(ca_n\pm db_n)=c\sum\limits_{n=1}^{\infty}+d\sum\limits_{n=1}^{\infty}b_n$为收敛的且$\lim\limits_{n\to\infty}a_n=\lim\limits_{n\to\infty}b_n=0$
\end{dingli}

\begin{dingli}
对实数域$\R$\\
若有级数$\sum\limits_{n=1}^{\infty}a_n$和$\sum\limits_{n=1}^{\infty}b_n$\\
则级数$\sum\limits_{n=1}^{\infty}a_n$为收敛的当且仅当\\
对任意$\varepsilon>0$,存在$N\in\N^*$,对任意$m>n\geq N$,有$|\sum\limits_{i=n}^{m}a_i|\leq \varepsilon$\\
若对任意$n\in\N^*$,有$a_n\geq b_n\geq 0$且$\sum\limits_{n=1}^{\infty}a_n$为收敛的,则$\sum\limits_{n=1}^{\infty}b_n$为收敛的\\
若$\sum\limits_{n=1}^{\infty}a_n$和$\sum\limits_{n=1}^{\infty}b_n$均为绝对收敛的,则$\sum\limits_{n=1}^{\infty}\sum\limits_{i=1}^{n}a_{i}b_{n+1-i}=\sum\limits_{n=1}^{\infty}a_n\sum\limits_{n=1}^{\infty}b_n$为绝对收敛的
\end{dingli}

\begin{dingli}
Dirichlet:\\
对实数域$\R$\\
若对级数$\sum\limits_{n=1}^{\infty}a_n$和$\sum\limits_{n=1}^{\infty}b_n$\\
有$\{a_n\}$单调,即$a_1\leq a_2\leq \cdots$或$a_1\geq a_2\geq \cdots$\\
有$\{a_n\}$收敛于$0$,即$\lim\limits_{n\to\infty}a_n=0$\\
有$\{b_n\}$部分和一致有界,即存在$M>0$,对任意$n\in\N^*$,有$|\sum\limits_{i=1}^{n}b_n|\leq M$\\
则$\sum\limits_{n=1}^{\infty}a_nb_n$为收敛的
\end{dingli}

\begin{dingli}
Abel:\\
对实数域$\R$\\
若对级数$\sum\limits_{n=1}^{\infty}a_n$和$\sum\limits_{n=1}^{\infty}b_n$\\
有$\{a_n\}$单调,即$a_1\leq a_2\leq \cdots$或$a_1\geq a_2\geq \cdots$\\
有$\{a_n\}$有界,即存在$M>0$,对任意$n\in \N^*$,有$|a_n|<M$\\
有$\{b_n\}$级数收敛,即$\sum\limits_{n=1}^{\infty}b_n<\infty$\\
则$\sum\limits_{n=1}^{\infty}a_nb_n$为收敛的
\end{dingli}

\begin{dingli}
\[\e=\lim\limits_{n\to\infty}(1+\frac{1}{n})^{n}=\sum\limits_{n=0}^{\infty}\frac{1}{n!}\]
\end{dingli}

\chapter{函数}

\begin{dingyi}
对实数域$\R$\\
则称映射$f:D\subset \R\to \R$为函数,则称$D$为$f$的定义域,则称$f(D)$为值域\\
若$f(D)$为有界的,则称$f(x)$在$D$上有界\\
若存在$\alpha,a\in \R$,对任意$\varepsilon>0$,存在$\delta>0$,对任意$0<a-x<\delta$,有$|f(x)-\alpha|<\varepsilon$\\
则称$x$趋于$a^{-}$时$f(x)$有左极限$\alpha$,记为$\lim\limits_{x\to a^{-}}f(x)=\alpha$\\
若存在$\alpha,a\in \R$,对任意$\varepsilon>0$,存在$\delta>0$,对任意$0<x-a<\delta$,有$|f(x)-\alpha|<\varepsilon$\\
则称$x$趋于$a^{+}$时$f(x)$有右极限$\alpha$,记为$\lim\limits_{x\to a^{+}}f(x)=\alpha$\\
若存在$\alpha,a\in \R$,对任意$\varepsilon>0$,存在$\delta>0$,对任意$0<|x-a|<\delta$,有$|f(x)-\alpha|<\varepsilon$\\
则称$x$趋于$a$时$f(x)$有极限$\alpha$,记为$\lim\limits_{x\to a}f(x)=\alpha$\\
若存在$a\in \R$,对任意$M\in \R$,存在$\delta>0$,对任意$0<a-x<\delta$,有$f(x)>M$\\
则称$x$趋于$a$时$f(x)$有极限$\infty$,记为$\lim\limits_{x\to a^{-}}f(x)=\infty$\\
若存在$a\in \R$,对任意$M\in \R$,存在$\delta>0$,对任意$0<x-a<\delta$,有$f(x)>M$\\
则称$x$趋于$a$时$f(x)$有极限$\infty$,记为$\lim\limits_{x\to a^{+}}f(x)=\infty$\\
若存在$a\in \R$,对任意$M\in \R$,存在$\delta>0$,对任意$0<|x-a|<\delta$,有$f(x)>M$\\
则称$x$趋于$a$时$f(x)$有极限$\infty$,记为$\lim\limits_{x\to a}f(x)=\infty$\\
若存在$a\in \R$,对任意$M\in \R$,存在$\delta>0$,对任意$0<a-x<\delta$,有$f(x)<M$\\
则称$x$趋于$a$时$f(x)$有极限$-\infty$,记为$\lim\limits_{x\to a^{-}}f(x)=-\infty$\\
若存在$a\in \R$,对任意$M\in \R$,存在$\delta>0$,对任意$0<x-a<\delta$,有$f(x)<M$\\
则称$x$趋于$a$时$f(x)$有极限$-\infty$,记为$\lim\limits_{x\to a^{+}}f(x)=-\infty$\\
若存在$a\in \R$,对任意$M\in \R$,存在$\delta>0$,对任意$0<|x-a|<\delta$,有$f(x)<M$\\
则称$x$趋于$a$时$f(x)$有极限$-\infty$,记为$\lim\limits_{x\to a}f(x)=-\infty$
\end{dingyi}

\begin{zhu}
若$x$趋于$a$时$f$有$($左/右$)$极限为$\infty$/$-\infty$,也可称$f$在$a$处发散
\end{zhu}

\begin{dingli}
对实数域$\R$\\
若对函数$f,g:D\subset \R\to \R$,有$\lim\limits_{x\to a}f(x)=\alpha$且$\lim\limits_{x\to a}g(x)=\beta$\\
则有$\lim\limits_{x\to a}f(x)+g(x)=\alpha+\beta$且$\lim\limits_{x\to a}f(x)g(x)=\alpha\beta$且$\lim\limits_{x\to a}\frac{f(x)}{g(x)}=\frac{\alpha}{\beta},(g(x)\neq0,\beta\neq0)$
\end{dingli}

\begin{dingli}
Heine:\\
对实数域$\R$\\
若对函数$f:D\subset \R\to \R$,则$\lim\limits_{x\to a}f(x)=\alpha$当且仅当\\
对任意数列$\{x_n\}\subset D$,有$\lim\limits_{n\to\infty}x_n=a$,则$\lim\limits_{n\to \infty}f(x_n)=\alpha$
\end{dingli}

\begin{dingyi}
对实数域$\R$\\
若对函数列$\{f_n(x)\}$和函数$g(x):D\subset \R\to \R$\\
对$a\in \R$,存在$\alpha\in \R$,有$\lim\limits_{n\to \infty}f_n(a)=\alpha$\\
则称$\{f_n(x)\}$在$x=a$处收敛于$\alpha$\\
若对任意$x_0\in D$,有$\lim\limits_{n\to\infty}f_n(x_0)=g(x_0)$\\
则称$\{f_n(x)\}$在$D$上逐点收敛于$g(x)$,记为$\lim\limits_{n\to\infty}f_n(x)=g(x)$\\
对任意$\varepsilon>0$,存在$N\in\N$,对任意$x\in D,n\geq N$,有$|f_n(x)-g(x)|<\varepsilon$\\
则称$\{f_n(x)\}$在$D$上一致收敛于$g(x)$\\
记$s_n(x)=\sum\limits_{i=1}^{n}f_i(x),\sum\limits_{n=1}^{\infty}f_n(x)=\lim\limits_{n\to\infty}s_n(x)$\\
则称$s_n(x)$为$\{f_n(x)\}$的部分和,则称$\sum\limits_{n=1}^{\infty}f_n(x)$为函数项级数\\
若$\{s_n(x)\}$在$D$上逐点收敛于$g(x)$\\
则称$\sum\limits_{n=1}^{\infty}f_n(x)$在$D$上逐点收敛于$g(x)$,记为$\sum\limits_{n=1}^{\infty}f_n(x)=g(x)$\\
若$\{s_n(x)\}$在$D$上一致收敛于$g(x)$\\
则称$\sum\limits_{n=1}^{\infty}f_n(x)$在$D$上一致收敛于$g(x)$
\end{dingyi}

\begin{dingli}
Cauchy:\\
对实数域$\R$\\
若对函数列$\{f_n(x)\}$和函数$g(x):D\subset \R\to \R$\\
则$\{f_n(x)\}$则$D$上一致收敛于$g(x)$当且仅当\\
对任意$\varepsilon>0$,存在$N\in\N$,对任意$x\in D,n,m\geq N$,有$|f_n(x)-f_m(x)|<\varepsilon$
\end{dingli}

\begin{dingli}
Weierstrass:\\
对实数域$\R$\\
若对函数列$\{f_n(x)\}$和正项级数$\sum\limits_{n=1}^{\infty}a_n<\infty$\\
若对任意$n\in\N^*$,对任意$x\in D$,有$|f_n(x)|\leq a_n$,则$\sum\limits_{n=1}^{\infty}f_n(x)$在$D$上一致收敛
\end{dingli}

\begin{dingli}
Dirichlet:\\
对实数域$\R$\\
若对函数列$\{f_n(x)\}$和$\{g_n(x)\}$\\
有$\{f_n(x)\}$在$D$上关于$n$单调不增或单调不减\\
有$\{f_n(x)\}$在$D$上一致收敛于$0$\\
有$\{g_n(x)\}$在$D$上部分和一致有界,即存在$M>0$,对任意$n\in \N^*$,有$|\sum\limits_{i=1}^{n}g_n(x)|\leq M$\\
则$\sum\limits_{n=1}^{\infty}f_n(x)g_n(x)$在$D$上一致收敛
\end{dingli}

\begin{dingli}
Abel:\\
对实数域$\R$\\
若对函数列$\{f_n(x)\}$和$\{g_n(x)\}$\\
有$\{f_n(x)\}$在$D$上关于$n$单调不增或单调不减\\
有$\{f_n(x)\}$在$D$上一致有界,即存在$M>0$,对任意$n\in \N^*,x\in D$,有$|g_n(x)|\leq M$\\
有$\sum\limits_{n=1}^{\infty}g_n(x)$在$D$上一致收敛\\
则$\sum\limits_{n=1}^{\infty}f_n(x)g_n(x)$在$D$上一致收敛 
\end{dingli}

\begin{dingyi}
对实数域$\R$\\
对$\{a_n\}\subset \R$,则称$\sum\limits_{n=0}^{\infty}a_nx^n$为幂级数\\
若存在$\rho>0$,若$|x|<\rho$,有$\sum\limits_{n=0}^{\infty}a_nx^n$为绝对收敛的;若$|x|>\rho$,有$\sum\limits_{n=0}^{\infty}a_nx^n$为发散的\\
则称$\rho$为$\sum\limits_{n=0}^{\infty}a_nx^n$的收敛半径
\end{dingyi}

\begin{dingli}
Cauchy-Hadamard:\\
对实数域$\R$\\
则幂级数$\sum\limits_{n=0}^{\infty}a_nx^n$有收敛半径$\rho=\frac{1}{\varlimsup\limits_{n\to\infty}|a_n|^{\frac{1}{n}}}$
\end{dingli}

\begin{zhu}
若$\varlimsup\limits_{n\to\infty}|a_n|^{\frac{1}{n}}=0$,则$\rho=\infty$;若$\varlimsup\limits_{n\to\infty}|a_n|^{\frac{1}{n}}=\infty$,则$\rho=0$
\end{zhu}

\begin{dingli}
d'Alembert:\\
对实数域$\R$\\
则幂级数$\sum\limits_{n=0}^{\infty}a_nx^n$有收敛半径$\rho=\lim\limits_{n\to\infty}|\frac{a_n}{a_{n+1}}|$
\end{dingli}

\begin{dingli}
Abel-1:\\
对实数域$\R$\\
若$x=x_0$,有$\sum\limits_{n=0}^{\infty}a_nx^n$为收敛的,则对任意$|x|<|x_0|$,有$\sum\limits_{n=0}^{\infty}a_nx^n$为绝对收敛的\\
若$x=x_0$,有$\sum\limits_{n=0}^{\infty}a_nx^n$为发散的,则对任意$|x|>|x_0|$,有$\sum\limits_{n=0}^{\infty}a_nx^n$为发散的
\end{dingli}

\begin{dingli}
Abel-2:\\
对实数域$\R$\\
若$x=x_0>0$,有$\sum\limits_{n=0}^{\infty}a_nx^n$为收敛的\\
则对任意闭区间$U\subset(-x_0,x_0]$,有$\sum\limits_{n=0}^{\infty}a_nx^n$为一致收敛的
\end{dingli}

\begin{dingyi}
对实数域$\R$\\
则称$\e^x=\sum\limits_{n=0}^{\infty}\frac{1}{n!}x^n$为指数函数\\
则称$\cos x=\sum\limits_{n=0}^{\infty}(-1)^{n}\frac{x^{2n}}{(2n)!}$为余弦函数,则称$\sec x=\frac{1}{\cos x}$为正割函数\\
则称$\sin x=\sum\limits_{n=0}^{\infty}(-1)^{n}\frac{x^{2n+1}}{(2n+1)!}$为正弦函数,则称$\csc x=\frac{1}{\sin x}$为余割函数\\
则称$\tan x=\frac{\sin x}{\cos x}$为正切函数,则称$\cot x=\frac{1}{\tan x}$为余切函数
\end{dingyi}

\begin{dingyi}
对复数域$\C$\\
则称$\e^x=\sum\limits_{n=0}^{\infty}\frac{1}{n!}x^n$为指数函数\\
则称$\cos x=\frac{\e^{\i x}+\e^{-\i x}}{2}$为余弦函数,则称$\sec x=\frac{1}{\cos x}$为正割函数\\
则称$\sin x=\frac{\e^{\i x}-\e^{-\i x}}{2}$为正弦函数,则称$\csc x=\frac{1}{\sin x}$为余割函数\\
则称$\tan x=\frac{\e^{\i x}-\e^{-\i x}}{\e^{\i x}+\e^{-\i x}}$为正切函数,则称$\cot x=\frac{1}{\tan x}$为余切函数
\end{dingyi}

\begin{zhu}
上述定义等价
\end{zhu}

\begin{dingli}
对复数域$\C$\\
对任意$x,y\in \C$,有$\e^{x}\e^{y}=\e^{x+y}$\\
有$\cos(x+y)=\cos x\cos y-\sin x\sin y$\\
有$\sin(x+y)=\sin x\cos y+\cos x\sin y$
\end{dingli}

\begin{dingyi}
Euler:\\
对复数域$\C$\\
对任意$x\in \C$,有$\e^{\i x}=\cos x+\i\sin x$
\end{dingyi}

\begin{dingli}
\[\lim\limits_{x\to 0}\e^x=1\]
\[\lim\limits_{x\to 0}\frac{\sin x}{x}=1\]
\end{dingli}

\begin{dingyi}
对实数域$\R$\\
对函数$f:D\subset \R \to \R$和$x_0\in D$\\
若有$\lim\limits_{x\to x_0^{-}}f(x)=f(x_0)$,则称$f(x)$在$x_0$处为左连续的\\
若有$\lim\limits_{x\to x_0^{+}}f(x)=f(x_0)$,则称$f(x)$在$x_0$处为右连续的\\
若有$\lim\limits_{x\to x_0}f(x)=f(x_0)$,则称$f(x)$在$x_0$处为连续的\\
即对任意$\varepsilon>0$,存在$\delta>0$,对任意$|x-x_0|<\delta$,有$|f(x)-f(x_0)|<\varepsilon$\\
若有$\lim\limits_{x\to x_0^{-}}f(x)<\infty$且$\lim\limits_{x\to x_0^{+}}f(x)<\infty$,但$\lim\limits_{x\to x_0^{-}}f(x)\neq f(x_0)$或$\lim\limits_{x\to x_0^{+}}f(x)\neq f(x_0)$\\
则称$x_0$为$f(x)$的第一类间断点\\
若有$|\lim\limits_{x\to x_0^{-}}f(x)|=\infty$或$|\lim\limits_{x\to x_0^{+}}f(x)|=\infty$或$\lim\limits_{x\to x_0^{-}}f(x)$不存在或$\lim\limits_{x\to x_0^{+}}f(x)$不存在\\
则称$x_0$为$f(x)$的第二类间断点
\end{dingyi}

\begin{dingli}
对实数域$\R$\\
对函数$f,g:D\subset \R \to \R$和$x_0\in D$\\
若$f(x),g(x)$均在$x_0$处为连续的\\
则$f(x)+g(x),f(x)-g(x),f(x)g(x),\frac{f(x)}{g(x)}(g(x_0)\neq 0)$在$x_0$处为连续的\\
若$f(x)$在$x_0$处为连续的且$g(x)$在$f(x_0)$处为连续的\\
则$(g\circ f)(x)=g(f(x))$在$x_0$处为连续的
\end{dingli}

\begin{dingli}
对实数域$\R$\\
对函数$f:[a,b]\subset \R \to \R$在任意$x_0\in [a,b]$处连续\\
若有$f(a)<0<f(b)$,则存在$c\in (a,b)$,有$f(c)=0$\\
若有$f(a)<C<f(b)$,则存在$c\in (a,b)$,有$f(c)=C$
\end{dingli}

\begin{dingli}
对实数域$\R$\\
对函数$f:[a,b]\subset \R \to \R$在任意$x_0\in [a,b]$处连续\\
则$f(x)$在$[a,b]$上有界\\
则存在$c_1,c_2\in [a,b]$,有$f(c_1)=\sup\limits_{x\in[a,b]}f(x),f(c_2)=\inf\limits_{x\in[a,b]}f(x)$
\end{dingli}

\begin{dingyi}
对实数域$\R$\\
对函数$f:D\subset \R \to \R$\\
若对任意$x_0\in D,\varepsilon>0$,存在$\delta>0$,对任意$|x-x_0|<\delta$,有$|f(x)-f(x_0)|<\varepsilon$\\
则称$f(x)$在$D$上连续\\
若对任意$\varepsilon>0$,存在$\delta>0$,对任意$x,y\in D,|x-y|<\delta$,有$|f(x)-f(y)|<\varepsilon$\\
则称$f(x)$在$D$上一致连续
\end{dingyi}

\begin{dingli}
对实数域$\R$\\
对函数$f:[a,b]\subset \R \to \R$\\
则$f(x)$在$[a,b]$上连续当且仅当$f(x)$在$[a,b]$上一致连续
\end{dingli}

\begin{zhu}
若$f(x)$在$x_0\in [a,b]$上连续\\
则对任意$\varepsilon>0$,存在$\delta>0$,对任意$x,y\in[x_0-\delta,x_0+\delta]\cap[a,b]$,有$|f(x)-f(y)|<\varepsilon$\\
对$f(x)$在$[a,c]\cup[c,d]\cup[d,b]$上连续\\
若对$\varepsilon>0$,存在$\delta_1>0,\delta_2>0$,有$f(x)$在$[a,d],[c,b]$上分别满足一致连续条件\\
则存在$\delta=\min\{\delta_1,\delta_2,d-c\}$有$f(x)$在$[a,b]$上满足一致连续条件
\end{zhu}

\begin{dingli}
对实数域$\R$\\
对函数列$\{f_n(x)\}$在$[a,b]$上一致收敛于$g(x)$\\
若对任意$n\in \N^*$,有$f_n(x)$在$[a,b]$上连续,则$g(x)$在$[a,b]$上连续\\
对函数项级数$\sum\limits_{n=1}^{\infty}f_n(x)$在$[a,b]$上一致收敛于$g(x)$\\
若对任意$n\in \N^*$,有$f_n(x)$在$[a,b]$上连续,则$g(x)$在$[a,b]$上连续
\end{dingli}

\begin{dingli}
Dini:\\
对实数域$\R$\\
对函数列$\{f_n(x)\}$在$[a,b]$上逐点收敛于$g(x)$\\
若对任意$n\in \N^*$,有$f_n(x)$在$[a,b]$上连续且对任意$x\in[a,b]$,有$f_n(x)\leq f_{n+1}(x)$\\
则$\{f_n(x)\}$在$[a,b]$上一致收敛于$g(x)$
\end{dingli}

\begin{zhu}
对任意$x_0\in [a,b],\varepsilon>0$,存在$\delta>0,N\in\N$\\
对任意$x\in[x_0-\delta,x_0+\delta],n\geq N$,有$|f_n(x)-g(x)|<\varepsilon$
\end{zhu}

\begin{dingyi}
对实数域$\R$\\
对函数$f:D\subset \R \to \R$\\
若对任意$x,y\in D,x<y$,有$f(x)<f(y)$,则称$f(x)$在$D$上单调递增\\
若对任意$x,y\in D,x<y$,有$f(x)>f(y)$,则称$f(x)$在$D$上单调递减
\end{dingyi}

\begin{dingli}
对实数域$\R$\\
对函数$f:[a,b]\subset \R \to \R$\\
若$f(x)$在$[a,b]$上连续且单调递增,则$f$为$[a,b]\to [f(a),f(b)]$的双射\\
则逆映射$f^{-1}$在$[f(a),f(b)]\to [a,b]$上连续且单调递增,称为$f(x)$的反函数\\
若$f(x)$在$[a,b]$上连续且单调递减,则$f$为$[a,b]$与$[f(b),f(a)]$的双射\\
则逆映射$f^{-1}$在$[f(a),f(b)]\to [a,b]$上连续且单调递减,称为$f(x)$的反函数
\end{dingli}

\begin{dingyi}
对实数域$\R$\\
则有$\e^x$在$\R\to(0,\infty)$上单调递增\\
则称$\ln x:(0,\infty)\to\R$为$(\e^x)$的反函数,为对数函数\\
则有$\cos x$在$[0,\pi]\to[-1,1]$上单调递减\\
则称$\arccos x:[-1,1]\to[0,\pi]$为$\cos x$的反函数,为反余弦函数\\
则有$\sin x$在$[-\frac{\pi}{2},\frac{\pi}{2}]\to[-1,1]$上单调递增\\
则称$\arcsin x:[-1,1]\to[-\frac{\pi}{2},\frac{\pi}{2}]$为$\sin x$的反函数,为反正弦函数\\
则有$\tan x$在$[-\frac{\pi}{2},\frac{\pi}{2}]\to\R$上单调递增\\
则称$\arctan x:[-\frac{\pi}{2},\frac{\pi}{2}]\to\R$为$\tan x$的反函数,为反正切函数
\end{dingyi}

\begin{dingyi}
对实数域$\R$\\
则称$x^a=\e^{a\ln x}:(0,\infty)\to\R$为幂函数\\
若$a>0,a\neq 1$,则有$\e^{x\ln a}$在$\R\to (0,\infty)$单调\\
则称$a^x=\e^{x\ln a}$为一般指数函数\\
则称$\log_{a}x$为$a^x$的反函数,为一般对数函数
\end{dingyi}

\begin{dingli}
对实数域$\R$\\
对任意$x>0,y,z\in \R$,有$x^{y+z}=x^yx^z$\\
对任意$x>0,y,z\in \R$,有$x^{yz}=(x^y)^z$\\
对任意$x>0,y\in \R$,有$\ln(x^{y})=y\ln x$\\
对任意$x,y>0$,有$\ln(x+y)=\ln x+\ln y$\\
对任意$x,y>0$,有$\log_{x}y=\frac{\ln y}{\ln x}$
\end{dingli}

\chapter{微分}

\begin{dingyi}
对实数域$\R$\\
对函数$f,g:(a,b)\subset \R \to \R$,其中$c\in (a,b)$\\
若存在$\varepsilon>0$,对任意$|x-c|<\varepsilon$,有$g(x)\neq 0$且$\lim\limits_{x\to c}\frac{f(x)}{g(x)}=0$\\
则称在$x\to c$时$f(x)$为$g(x)$的高阶无穷小\\
记为$f(x)=o(g(x)),x\to c$\\
若存在$\varepsilon>0$,对任意$|x-c|<\varepsilon$,有$g(x)\neq 0$且$|\lim\limits_{x\to c}\frac{f(x)}{g(x)}|<\infty$\\
则称在$x\to c$时$f(x)$为$g(x)$的非低阶无穷小\\
记为$f(x)=O(g(x)),x\to c$
\end{dingyi}

\begin{zhu}
以上$c$还可取$c^{-},c^{+},\infty,-\infty$
\end{zhu}

\begin{dingyi}
对实数域$\R$\\
对函数$f:[a,b]\subset \R \to \R$,其中$x_0\in [a,b]$\\
若极限$\lim\limits_{h\to 0}\frac{f(x+h)-f(x)}{h}$存在,记$f'(x_0)=\lim\limits_{h\to 0}\frac{f(x_0+h)-f(x_0)}{h}$\\
则称$f(x)$在$x_0$处可导,则称$f'(x_0)$为$f(x)$在$x_0$处的导数\\
则称$f'_{-}(x_0)=\lim\limits_{h\to 0^{-}}\frac{f(x_0+h)-f(x_0)}{h}$为$f(x)$在$x_0$处的左导数\\
则称$f'_{+}(x_0)=\lim\limits_{h\to 0^{+}}\frac{f(x_0+h)-f(x_0)}{h}$为$f(x)$在$x_0$处的右导数
\end{dingyi}

\begin{dingyi}
对实数域$\R$\\
对函数$f:[a,b]\subset \R \to \R$,其中$x_0\in [a,b]$\\
记$\Delta f(x_0)=f(x_0+h)-f(x_0),\Delta x=(x_0+h)-x_0$\\
若在$\Delta x\to 0$时,存在只与$x_0$有关的数$\alpha_{x_0}\in\R$,有$\Delta f(x_0)=\alpha_{x_0} \Delta x+o(\Delta x)$\\
则称$f(x)$在$x_0$处可微,则称$\d f(x_0)=\alpha \d x$为$f(x)$在$x_0$处的微分\\
其中$\d x=\Delta x$为$x$在$x_0$处的微分\\
则记$\frac{\d f(x)}{\d x}|_{x=x_0}=\alpha$,则称$\frac{\d f(x)}{\d x}:[a,b]\to \R$为微分算子
\end{dingyi}

\begin{dingli}
对实数域$\R$\\
对函数$f:[a,b]\subset \R \to \R$,其中$x_0\in [a,b]$\\
则$f(x)$在$x_0$处可导当且仅当$f(x)$在$x_0$处可微\\
有$\frac{\d f(x)}{\d x}|_{x=x_0}=f'(x_0)$且$f(x)$在$x_0$处连续
\end{dingli}

\begin{dingli}
对实数域$\R$\\
对函数$f,g:[a,b]\subset \R \to \R$,有$f(x)$和$g(x)$在$x_0\in [a,b]$处可导\\
则对任意$\alpha,\beta\in\R$,有$(\alpha f+\beta g)'(x_0)=\alpha f'(x_0)+\beta g'(x_0)$\\
则$(fg)'(x_0)=f'(x_0)g(x_0)+f(x_0)g'(x_0)$\\
则$(\frac{f}{g})'(x_0)=\frac{f'(x_0)g(x_0)-f(x_0)g'(x_0)}{(g(x_0))^2},g(x_0)\neq 0$\\
若$g(x)$在$f(x_0)\in [a,b]$处可导,则$(g\circ f)'(x_0)=g'(f(x_0))f'(x_0)$\\
若$f(x)$在$[a,b]$上连续且单调,则$(f^{-1})'(f(x_0))=\frac{1}{f'(x_0)}$
\end{dingli}

\begin{dingli}
对实数域$\R$\\
对函数$f,g:[a,b]\subset \R \to \R$,有$f(x)$和$g(x)$在$x_0\in [a,b]$处可微\\
则对任意$\alpha,\beta\in\R$,有$\d(\alpha f+\beta g)(x_0)=\alpha \d f(x_0)+\beta \d g(x_0)$\\
则$\d(fg)(x_0)=\d f(x_0)g(x_0)+f(x_0)\d g(x_0)$\\
则$\d(\frac{f}{g})(x_0)=\frac{\d f(x_0)g(x_0)-f(x_0)\d g(x_0)}{g(x_0)^2}$\\
若$g(x)$在在$f(x_0)\in [a,b]$处可微,则$\d (g\circ f)(x_0)=f'(g(x_0))\d g(x_0)$\\
若$f(x)$在$[a,b]$上连续且单调,则$\d f^{-1}(f(x_0))=\frac{1}{f'(x_0)}\d f(x_0)$
\end{dingli}

\begin{dingyi}
对实数域$\R$\\
对函数$f:[a,b]\subset \R \to \R$\\
若对任意$x_0\in (a,b)$,有$f(x)$在$x_0$处可导或可微\\
则称$f(x)$在$(a,b)$上可导或可微\\
若对任意$x_0\in (a,b)$,有$f(x)$在$x_0$处可导或可微且$f'_{+}(a)$和$f'_{+}(b)$存在\\
则称$f(x)$在$[a,b]$上可导或可微\\
则称$f':[a,b]\to \R,x_0\to f'(x_0)$为$f(x)$在$[a,b]$上的导函数\\
则称$\d f:[a,b]\to \R\d x,x_0\to f'(x_0)\d x$为$f(x)$在$[a,b]$上的微分形式\\
记$f=f^{(0)},f'=f^{(1)}=(f^{(0)})',\cdots,f^{(n)}=(f^{(n-1)})',\cdots$\\
则称$f^{(n)}:[a,b]\to \R$为$f(x)$在$[a,b]$上的$n$阶导函数\\
若$f(x)$在$[a,b]$上存在$n$阶导函数,则称$f(x)$在$[a,b]$上的$n$阶可导\\
记$\d x=\d^{0} f,\d f=\d f^{1},\cdots,\d^{n}f=\d(\d^{n-1}f),\cdots$,其中$\d(\d x)=0$\\
则称$\d^{n}f:[a,b]\to \R\d x$为$f(x)$在$[a,b]$上的$n$阶微分形式\\
若$f(x)$在$[a,b]$上存在$n$阶微分形式,则称$f(x)$在$[a,b]$上的$n$阶可微
\end{dingyi}

\begin{zhu}
若$f(x)$在$[a,b]$上$n$阶可导或可微且$n$阶导函数连续,则称$n$阶连续可导或可微
\end{zhu}

\begin{dingli}
对实数域$\R$\\
对函数$f:[a,b]\subset \R \to \R$\\
则$f(x)$在$[a,b]$上可导当且仅当$f(x)$在$[a,b]$上可微\\
有$\frac{\d f}{\d x}=f'(x)$且$f(x)$在$[a,b]$上连续
\end{dingli}

\begin{dingli}
对实数域$\R$\\
对函数$f,g:[a,b]\subset \R \to \R$,有$f(x)$和$g(x)$在$[a,b]$上可导\\
则对任意$\alpha,\beta\in\R$,有$(\alpha f+\beta g)'=\alpha f'+\beta g'$\\
则$(fg)'=f'g+fg'$\\
则$(\frac{f}{g})'=\frac{f'g-fg'}{g^2},g(x)\neq 0$\\
若$g(x)$在$f([a,b])$上可导,则$(g\circ f)'(x)=g'(f(x))f'(x)$\\
若$f(x)$在$[a,b]$上连续且单调,则$(f^{-1})'(f(x))=\frac{1}{f'(x)}$
\end{dingli}

\begin{dingli}
对实数域$\R$\\
对函数$f,g:[a,b]\subset \R \to \R$,有$f(x)$和$g(x)$在$[a,b]$上可微\\
则对任意$\alpha,\beta\in\R$,有$\d(\alpha f+\beta g)=\alpha \d f+\beta \d g$\\
则$\d(fg)=\d fg+f\d g$\\
则$\d(\frac{f}{g})=\frac{\d fg-f\d g}{g^2},g\neq 0$\\
若$g(x)$在$f([a,b])$上可微,则$\d (g\circ f)(x)=f'(g(x))\d g(x)$\\
若$f(x)$在$[a,b]$上连续且单调,则$\d f^{-1}(f(x))=\frac{1}{f'(x)}\d f(x)$
\end{dingli}

\begin{dingli}
Leibniz:\\
对实数域$\R$\\
对函数$f,g:[a,b]\subset \R \to \R$,有$f(x)$和$g(x)$在$[a,b]$上可导或可微\\
则$(f+g)^{(n)}=f^{(n)}+g^{(n)}$且$\d^{n}(f+g)=\d^{n}f+\d^{n}g$\\
则$(fg)^{(n)}=\sum\limits_{i=0}^{n}C_{n}^{i}f^{(i)}g^{(n-i)}$且$\d^{n}(fg)=\sum\limits_{i=0}^{n}C_{n}^{i}(\d^{i}f\otimes\d^{n-i}g)$
\end{dingli}

\begin{zhu}
以下为初等函数的导函数:\\
\[(C)'=0,x^{\alpha}=\alpha x^{\alpha-1}\]
\[(\e^{x})'=\e^{x},(\ln x)'=\frac{1}{x}\]
\[(a^x)'=a^x\ln a,(\log_{a}x)'=\frac{1}{x}\log_{a}\e\]
\[(\cos x)'=-\sin x,(\sin x)'=\cos x,(\tan x)'=\sec^2 x\]
\[(\sec x)'=\tan x\sec x,(\csc x)'=-\cot x\csc x,(\cot x)'=-\csc^2 x\]
\[(\arccos x)'=-\frac{1}{\sqrt{1-x^2}},(\arcsin x)'=\frac{1}{\sqrt{1-x^2}},(\arctan x)'=\frac{1}{1+x^2}\]
\end{zhu}

\begin{dingyi}
对实数域$\R$\\
对函数$f:[a,b]\subset \R \to \R$,其中$c\in(a.b)$\\
若存在$\varepsilon>0$,对任意$|x-c|<\varepsilon$,有$f(x)\leq f(c)$\\
则称$c$为$f(x)$的极大点,则称$f(c)$为极大值\\
若存在$\varepsilon>0$,对任意$|x-c|<\varepsilon$,有$f(x)\geq f(c)$\\
则称$c$为$f(x)$的极小点,则称$f(c)$为极小值\\
若对任意$x\in[a,b]$,有$f(x)\leq f(c)$,则称$c$为$f(x)$的最大点,则称$f(c)$为最大值\\
若对任意$x\in[a,b]$,有$f(x)\geq f(c)$,则称$c$为$f(x)$的最小值,则称$f(c)$为最小值\\
若有$f(x)$在$[a,b]$上可导或可微且$f'(c)=0$\\
则称$c$为$f(x)$的临界点
\end{dingyi}

\begin{dingli}
Fermat:\\
对实数域$\R$\\
对函数$f:[a,b]\subset \R \to \R$,有$f(x)$在$[a,b]$上可导或可微\\
若$c\in (a,b)$为$f(x)$的极大点或极小点,则$c$为$f(x)$的临界点
\end{dingli}

\begin{dingli}
Rolle:\\
对实数域$\R$\\
对函数$f:[a,b]\subset \R \to \R$,有$f(x)$在$(a,b)$上可导或可微且在$[a,b]$上连续\\
若$f(a)=f(b)$,则存在$c\in (a,b)$,有$c$为$f(x)$的临界点
\end{dingli}

\begin{dingli}
Lagrange:\\
对实数域$\R$\\
对函数$f:[a,b]\subset \R \to \R$,有$f(x)$在$(a,b)$上可导或可微且在$[a,b]$上连续\\
则存在$c\in [a,b]$,有$f(b)-f(a)=f'(c)(b-a)$
\end{dingli}

\begin{dingli}
Darboux:\\
对实数域$\R$\\
对函数$f:[a,b]\subset \R \to \R$,有$f(x)$在$[a,b]$上可导或可微\\
若$f'(a)<C<f'(b)$或$f'(a)>C>f'(b)$,则存在$c\in (a,b)$,有$f'(c)=C$
\end{dingli}

\begin{dingli}
Cauchy:\\
对实数域$\R$\\
对函数$f,g:[a,b]\subset \R \to \R$,有$f(x)$和$g(x)$在$[a,b]$上可导或可微且在$[a,b]$上连续\\
若对任意$x\in (a,b)$,有$g'(x)\neq 0$\\
则存在$c\in [a,b]$,有$\frac{f(b)-f(a)}{g(b)-g(a)}=\frac{f'(c)}{g'(c)}$
\end{dingli}

\begin{dingli}
l'Hôpital:\\
对实数域$\R$\\
对函数$f,g:(a,b)\subset \R \to \R$,其中$-\infty\leq a<b\leq \infty$\\
有$f(x)$和$g(x)$在$(a,b)$上$n$次可导或可微且对任意$x\in (a,b)$,有$g^{(n)}(x)\neq 0$\\
若有$-\infty\leq\lim\limits_{x\to a^{+}}\frac{f^{(n)}(x)}{g^{(n)}(x)}\leq \infty$\\
且对任意$i=0,1,\cdots,n-1$,有$\lim\limits_{x\to a^{-}}f^{(i)}(x)=\lim\limits_{x\to a^{+}}g^{(i)}(x)=0$或$\infty$\\
则有$\lim\limits_{x\to a^{+}}\frac{f(x)}{g(x)}=\lim\limits_{x\to a^{+}}\frac{f^{(n)}(x)}{g^{(n)}(x)}$\\
若有$-\infty\leq\lim\limits_{x\to b^{-}}\frac{f^{(n)}(x)}{g^{(n)}(x)}\leq \infty$\\
且对任意$i=0,1,\cdots,n-1$,有$\lim\limits_{x\to a^{-}}f^{(i)}(x)=\lim\limits_{x\to b^{-}}g^{(i)}(x)=0$或$\infty$\\
则有$\lim\limits_{x\to b^{-}}\frac{f(x)}{g(x)}=\lim\limits_{x\to b^{-}}\frac{f^{(n)}(x)}{g^{(n)}(x)}$
\end{dingli}

\begin{dingli}
Peano:\\
对实数域$\R$\\
对函数$f:[a,b]\subset \R \to \R$,其中$x_0\in (a,b)$\\
若$f(x)$在$[a,b]$上$n$阶可导或可微\\
则对任意$x\in [a,b]$,有$f(x)=f(x_0)+\frac{f'(x_0)}{1!}(x-x_0)+\cdots+\frac{f^{(n)}(x_0)}{n!}(x-x_0)^{n}+P_n(x)$\\
其中$P_n(x)=o((x-x_0)^n)$
\end{dingli}

\begin{dingli}
Taylor:\\
对实数域$\R$\\
对函数$f:[a,b]\subset \R \to \R$,其中$x_0\in (a,b)$\\
若$f(x)$在$[a,b]$上$n$阶连续可导或可微且在$(a,b)$上$n+1$阶可导或可微\\
则对任意$x\in [a,b]$,有$f(x)=f(x_0)+\frac{f'(x_0)}{1!}(x-x_0)+\cdots+\frac{f^{(n)}(x_0)}{n!}(x-x_0)^{n}+T_n(x)$\\
其中$T_n(x)=\frac{1}{n!}\frac{g(x)-g(x_0)}{g'(\xi)}f^{(n+1)}(\xi)(x-\xi)^{n}$且$\xi\in [a,b]$
\end{dingli}

\begin{dingli}
Roche:\\
对实数域$\R$\\
对函数$f:[a,b]\subset \R \to \R$,其中$x_0\in (a,b)$\\
若$f(x)$在$[a,b]$上$n$阶连续可导或可微且在$(a,b)$上$n+1$阶可导或可微\\
则对任意$x\in [a,b]$,有$f(x)=f(x_0)+\frac{f'(x_0)}{1!}(x-x_0)+\cdots+\frac{f^{(n)}(x_0)}{n!}(x-x_0)^{n}+R_n(x)$\\
其中$R_n(x)=\frac{1}{n!}\frac{(1-\theta)^{n+1-p}}{p}f^{(n+1)}(x_0+\theta(x-x_0))(x-x_0)^{n+1}$且$p>0,\theta\in (0,1)$
\end{dingli}

\begin{dingli}
Lagrange:\\
对实数域$\R$\\
对函数$f:[a,b]\subset \R \to \R$,其中$x_0\in (a,b)$\\
若$f(x)$在$[a,b]$上$n$阶连续可导或可微且在$(a,b)$上$n+1$阶可导或可微\\
则对任意$x\in [a,b]$,有$f(x)=f(x_0)+\frac{f'(x_0)}{1!}(x-x_0)+\cdots+\frac{f^{(n)}(x_0)}{n!}(x-x_0)^{n}+L_n(x)$\\
其中$L_n(x)=\frac{1}{(n+1)!}f^{(n+1)}(x_0+\theta(x-x_0))(x-x_0)^{n+1}$且$\theta\in (0,1)$
\end{dingli}

\begin{dingli}
Cauchy:\\
对实数域$\R$\\
对函数$f:[a,b]\subset \R \to \R$,其中$x_0\in (a,b)$\\
若$f(x)$在$[a,b]$上$n$阶连续可导或可微且在$(a,b)$上$n+1$阶可导或可微\\
则对任意$x\in [a,b]$,有$f(x)=f(x_0)+\frac{f'(x_0)}{1!}(x-x_0)+\cdots+\frac{f^{(n)}(x_0)}{n!}(x-x_0)^{n}+C_n(x)$\\
其中$C_n(x)=\frac{1}{n!}(1-\theta)^{n}f^{(n+1)}(x_0+\theta(x-x_0))(x-x_0)^{n+1}$且$\theta\in (0,1)$
\end{dingli}

\begin{dingyi}
对实数域$\R$\\
对函数$f:[a,b]\subset \R \to \R$,其中$x_0\in (a,b)$\\
若对任意$x\in [a,b]$,有$\lim\limits_{n\to \infty}T_n(x)=\lim\limits_{n\to\infty}(f(x)-\sum\limits_{i=0}^{n}\frac{f^{(i)}(x_0)}{n!}(x-x_0)^{i})=0$\\
即$T_n(x)$在$[a,b]$上逐点收敛于$0$\\
则$f(x)$在$[a,b]$上有Taylor展式,即$f(x)=\sum\limits_{n=0}^{\infty}\frac{f^{(n)}(x_0)}{n!}(x-x_0)^{n}$\\
若$x_0=0$,则称为Maclaurin展式,即$f(x)=\sum\limits_{n=0}^{\infty}\frac{f^{(n)}(0)}{n!}x^n$\\
若$T_n(x)$在$[a,b]$上一致收敛于$0$,则Taylor展式一致收敛于$f(x)$
\end{dingyi}

\begin{zhu}
以下为初等函数的Taylor展式:\\
\[\e^x=\sum\limits_{n=0}^{\infty}\frac{x^n}{n!}=1+x+\frac{x^2}{2}+\cdots,\  x\in\R\]
\[\cos x=\sum\limits_{n=0}^{\infty}(-1)^{n}\frac{x^{2n}}{(2n)!}=1-\frac{x^2}{2}+\frac{x^4}{24}+\cdots,\  x\in\R\]
\[\sin x=\sum\limits_{n=0}^{\infty}(-1)^{n+1}\frac{x^{2n+1}}{(2n+1)!}=x-\frac{x^3}{6}+\frac{x^5}{120}+\cdots,\  x\in\R\]
\[\ln(1+x)=\sum\limits_{n=0}^{\infty}(-1)^{n+1}\frac{x^{n}}{n}=x-\frac{x^2}{2}+\frac{x^3}{3}+\cdots,\  x\in(-1,1]\]
\[(1+x)^{\alpha}=\sum\limits_{n=0}^{\infty}\frac{\alpha(\alpha-1)\cdots(\alpha-n+1)}{n!}x^n=1+\alpha x+\frac{\alpha(\alpha-1)}{2}x^2+\cdots,\ x\in(-1,1),\alpha\in\R\]
\end{zhu}

\begin{dingyi}
对实数域$\R$\\
对函数$f:(a,b)\subset \R \to \R$\\
若对任意$x,y\in(a,b),\lambda\in(0,1)$,有$f(\lambda x+(1-\lambda)y)\leq \lambda f(x)+(1-\lambda)f(y)$\\
则称$f(x)$为$(a,b)$上的下凸函数\\
若对任意$x,y\in(a,b),\lambda\in(0,1)$,有$f(\lambda x+(1-\lambda)y)\geq \lambda f(x)+(1-\lambda)f(y)$\\
则称$f(x)$为$(a,b)$上的上凸函数\\
若对任意$x,y\in(a,b),\lambda\in(0,1)$,有$f(\lambda x+(1-\lambda)y)< \lambda f(x)+(1-\lambda)f(y)$\\
则称$f(x)$为$(a,b)$上的严格下凸函数\\
若对任意$x,y\in(a,b),\lambda\in(0,1)$,有$f(\lambda x+(1-\lambda)y)> \lambda f(x)+(1-\lambda)f(y)$\\
则称$f(x)$为$(a,b)$上的严格上凸函数
\end{dingyi}

\begin{dingli}
对实数域$\R$\\
对函数$f:(a,b)\subset \R \to \R$\\
则$f(x)$为$(a,b)$上的下凸函数当且仅当对任意$a<c<x<d<b$\\
有$\frac{f(d)-f(c)}{d-c}\leq \frac{f(d)-f(x)}{d-x}$或$\frac{f(c)-f(x)}{c-x}\leq \frac{f(d)-f(c)}{d-c}$或$\frac{f(c)-f(x)}{c-x}\leq \frac{f(d)-f(x)}{d-x}$\\
则$f(x)$为$(a,b)$上的上函数当且仅当对任意$a<c<x<d<b$\\
有$\frac{f(d)-f(c)}{d-c}\geq \frac{f(d)-f(x)}{d-x}$或$\frac{f(c)-f(x)}{c-x}\geq \frac{f(d)-f(c)}{d-c}$或$\frac{f(c)-f(x)}{c-x}\geq \frac{f(d)-f(x)}{d-x}$\\
则$f(x)$为$(a,b)$上的严格下凸函数当且仅当对任意$a<c<x<d<b$\\
有$\frac{f(d)-f(c)}{d-c}<\frac{f(d)-f(x)}{d-x}$或$\frac{f(c)-f(x)}{c-x}<\frac{f(d)-f(c)}{d-c}$或$\frac{f(c)-f(x)}{c-x}< \frac{f(d)-f(x)}{d-x}$\\
则$f(x)$为$(a,b)$上的严格上凸函数当且仅当对任意$a<c<x<d<b$\\
有$\frac{f(d)-f(c)}{d-c}>\frac{f(d)-f(x)}{d-x}$或$\frac{f(c)-f(x)}{c-x}>\frac{f(d)-f(c)}{d-c}$或$\frac{f(c)-f(x)}{c-x}> \frac{f(d)-f(x)}{d-x}$
\end{dingli}

\begin{dingli}
对实数域$\R$\\
对函数$f:(a,b)\subset \R \to \R$为$(a,b)$上的凸函数\\
则对任意$x_0\in(a,b)$,有$f'_{-}(x_0)$或$f'_{+}(x_0)$存在且$f(x)$在$x_0$处连续\\
则$\{x_0\in(a,b)|f'_{-}(x_0)\neq f'_{+}(x_0)\}$为至多可数集,即$f(x)$在$(a,b)$上不可导点至多可数
\end{dingli}

\begin{dingli}
对实数域$\R$\\
对函数$f:(a,b)\subset \R \to \R$在$(a,b)$上可导或可微\\
则$f(x)$为$(a,b)$上的下凸函数当且仅当对任意$a<c<d<b$,有$f'(c)\leq f'(d)$\\
则$f(x)$为$(a,b)$上的上函数当且仅当对任意$a<c<d<b$,有$f'(c)\geq f'(d)$\\
则$f(x)$为$(a,b)$上的严格下凸函数当且仅当对任意$a<c<d<b$,有$f'(c)< f'(d)$\\
则$f(x)$为$(a,b)$上的严格上凸函数当且仅当对任意$a<c<d<b$,有$f'(c)> f'(d)$
\end{dingli}

\begin{dingli}
对实数域$\R$\\
对函数$f:(a,b)\subset \R \to \R$在$(a,b)$上$2$阶可导或可微\\
则$f(x)$为$(a,b)$上的下凸函数当且仅当对任意$a<x<b$,有$f'(x)\geq 0$\\
则$f(x)$为$(a,b)$上的上函数当且仅当对任意$a<x<b$,有$f'(x)\leq 0$\\
对任意$a<x<b$,有$f'(x)> 0$,则$f(x)$为$(a,b)$上的严格下凸函数\\
对任意$a<x<b$,有$f'(x)< 0$,则$f(x)$为$(a,b)$上的严格上凸函数
\end{dingli}

\begin{dingli}
对实数域$\R$\\
对函数$f:(a,b)\subset \R \to \R$在$(a,b)$上可导或可微\\
若$f(x)$为$(a,b)$上的下凸函数且有临界点$c$,则$c$为最小点\\
若$f(x)$为$(a,b)$上的上凸函数且有临界点$c$,则$c$为最大点\\
若$f(x)$为$(a,b)$上的严格下凸函数且有临界点$c$,则$c$为唯一最小点\\
若$f(x)$为$(a,b)$上的严格上凸函数且有临界点$c$,则$c$为唯一最大点
\end{dingli}

\begin{dingli}
Jensen:\\
对实数域$\R$\\
对函数$f:(a,b)\subset \R \to \R$为$(a,b)$上的下凸函数\\
则对任意$x_1,\cdots,x_n\in(a,b),\lambda_1,\cdots,\lambda_n\in(0,1)$且$\lambda_1+\cdots+\lambda_n=1$\\
则有$f(\lambda_1x_1+\cdots+\lambda_n x_n)\leq \lambda_1f(x_1)+\cdots+\lambda_nf(x_n)$
\end{dingli}

\begin{dingli}
Young:\\
对实数域$\R$\\
对任意$a,b\in (0,\infty),p,q\in\R$且$p\neq 1,q\neq 1,\frac{1}{p}+\frac{1}{q}=1$\\
若$p>1$,则有$a^{\frac{1}{p}}b^{\frac{1}{q}}\leq \frac{a}{p}+\frac{b}{q}$;若$p<1$,则有$a^{\frac{1}{p}}b^{\frac{1}{q}}\geq \frac{a}{p}+\frac{b}{q}$\\
等号成立当且仅当$a=b$
\end{dingli}

\begin{dingli}
Hölder:\\
对实数域$\R$\\
对任意$x_1,\cdots,x_n,y_1,\cdots,y_n\in[0,\infty),p,q\in \R$且$p\neq 1,q\neq 1,\frac{1}{p}+\frac{1}{q}=1$\\
若$p>1$,则有$\sum\limits_{i=1}^{n}x_iy_i\leq (\sum\limits_{i=1}^{n}x_i^{p})^{\frac{1}{p}}(\sum\limits_{i=1}^{n}y_i^{q})^{\frac{1}{q}}$;若$p<1$,则有$\sum\limits_{i=1}^{n}x_iy_i\geq (\sum\limits_{i=1}^{n}x_i^{p})^{\frac{1}{p}}(\sum\limits_{i=1}^{n}y_i^{q})^{\frac{1}{q}}$\\
等号成立当且仅当存在$c\in \R$,对任意$i=1,\cdots n$,有$x_{i}^{p}=cy_{i}^{q}$
\end{dingli}

\begin{dingli}
Minkowski:\\
对实数域$\R$\\
对任意$x_1,\cdots,x_n,y_1,\cdots,y_n\in[0,\infty),p,q\in \R$且$p\neq 1,q\neq 1,\frac{1}{p}+\frac{1}{q}=1$\\
若$p>1$,则有$(\sum\limits_{i=1}^{n}(x_i+y_i)^{p})^{\frac{1}{p}}\leq (\sum\limits_{i=1}^{n}x_i^{p})^{\frac{1}{p}}+(\sum\limits_{i=1}^{n}y_i^{p})^{\frac{1}{p}}$;\\
若$p<1$,则有$(\sum\limits_{i=1}^{n}(x_i+y_i)^{p})^{\frac{1}{p}}\geq (\sum\limits_{i=1}^{n}x_i^{p})^{\frac{1}{p}}+(\sum\limits_{i=1}^{n}y_i^{p})^{\frac{1}{p}}$\\
等号成立当且仅当存在$c\in \R$,对任意$i=1,\cdots n$,有$x_{i}^{p}=cy_{i}^{q}$
\end{dingli}

\chapter{积分}

\begin{dingyi}
对实数域$\R$\\
对函数$f:[a,b]\subset \R \to \R$\\
则称$a=x_0<x_1<\cdots<x_n=b$为$[a,b]$的一个分划\\
若有$\xi_{i}\in [x_{i-1},x_{i}],i=1,2,\cdots,n$,记$\lambda=\max\limits_{i}\{x_i-x_{i-1}\}$\\
则称$\sum\limits_{i=1}^{n}f(\xi_{i})(x_{i}-x_{i-1})$为$f(x)$关于此分划的Riemann和\\
若$\lim\limits_{\lambda\to 0}\sum\limits_{i=1}^{n}f(\xi_{i})(x_{i}-x_{i-1})$存在且无关$\xi_{i}$的选取\\
则称$\int_{a}^{b}f(x)\d x=\lim\limits_{\lambda\to 0}\sum\limits_{i=1}^{n}f(\xi_{i})(x_{i}-x_{i-1})$为$f(x)$在$[a,b]$上的Riemann积分\\
则称$f(x)$在$[a,b]$上Riemann可积
\end{dingyi}

\begin{zhu}
此处极限的定义可表达为:对任意$\epsilon>0$,存在$\delta>0,I\in\R$,对任意$[a,b]$的分划\\
对任意$\xi_i\in[x_{i-1},x_i]$和$\max\limits_{i}\{x_i-x_{i-1}\}<\delta$,有$|I-\sum\limits_{i=1}^{n}f(\xi_{i})(x_{i}-x_{i-1})|<\varepsilon$
\end{zhu}

\begin{dingli}
对实数域$\R$\\
对函数$f:[a,b]\subset \R \to \R$\\
则$f(x)$在$[a,b]$上Riemann可积当且仅当\\
若有$[a,b]$上的分划$\{x_0,\cdots,x_n\}\subset\{y_0,\cdots,y_m\}$\\
对任意$\xi_i\in[x_{i-1},x_i],\eta_j\in[y_{j-1},y_j]$,若$\max\limits_{i}\{x_i-x_{i-1}\}<\delta$\\
则$|\sum\limits_{i=1}^{n}f(\xi_i)(x_i-x_{i-1})-\sum\limits_{j=1}^{m}f(\eta_{j})(y_j-y_{j-1})|<\varepsilon$
\end{dingli}

\begin{dingyi}
对实数域$\R$\\
对函数$f:[a,b]\subset \R \to \R$,其中$x_0\in(a,b)$\\
若有$[a,b]$上的分划$\{x_0,\cdots,x_n\}$\\
则称$\omega_{i}=\sup\limits_{\xi_{i}\in[x_{i-1},x_i]}f(\xi_{i})-\inf\limits_{\xi_{i}\in[x_{i-1},x_i]}f(\xi_i)$为$f(x)$在$[x_{i-1},x_i]$上的振幅\\
则称$\omega_{f}(x_0)=\lim\limits_{\varepsilon\to 0^{+}}(\sup\limits_{|x-x_0|<\varepsilon}f(x)-\inf\limits_{|x-x_0|<\varepsilon}f(x))$为$f(x)$在$x_0$处的振幅
\end{dingyi}

\begin{dingli}
对实数域$\R$\\
对函数$f:[a,b]\subset \R \to \R$\\
则$f(x)$在$[a,b]$上Riemann可积当且仅当\\
对任意$\varepsilon>0$,存在$\delta>0$,对任意$[a,b]$上的分划$\{x_0,\cdots,x_n\}$\\
若$\max\limits_{i}\{x_i-x_{i-1}\}<\delta$,则$\sum\limits_{i=1}^{n}\omega_{i}(x_i-x_{i-1})<\varepsilon$
\end{dingli}

\begin{dingli}
对实数域$\R$\\
对函数$f:[a,b]\subset \R \to \R$\\
若$f(x)$在$[a,b]$上Riemann可积,则$f(x)$在$[a,b]$上有界
\end{dingli}

\begin{dingli}
对实数域$\R$\\
对函数$f:[a,b]\subset \R \to \R$,其中$x_0\in(a,b)$\\
则$f(x)$在$x_0$处连续当且仅当$\omega_{f}(x_0)=0$\\
则$f(x)$在$[a,b]$上Riemann可积当且仅当对$D=\{x\in[a,b]|\omega_{f}(x)>0\}$\\
对任意$\varepsilon>0$,存在$\{(a_j,b_j)\}$,有$D\subset\bigcup\limits_{j=1}^{\infty}(a_j.b_j)$且$\sum\limits_{j=1}^{\infty}(b_j-a_j)<\varepsilon$\\
则若$f(x)$在$[a,b]$上连续,则$f(x)$在$[a,b]$上Riemann可积
\end{dingli}

\begin{dingli}
对实数域$\R$\\
对函数$f,g:[a,b]\subset \R \to \R$,其中$c\in(a,b)$\\
若$f(x)$在$[a,b]$上Riemann可积,则$f(x)$在$[a,c]$和$[c,b]$上Riemann可积\\
若$f(x)$在$[a,c]$和$[c,b]$上Riemann可积,则$f(x)$在$[a,b]$上Riemann可积\\
且$\int_{a}^{b}f(x)\d x=\int_{a}^{c}f(x)\d x+\int_{c}^{b}f(x)\d x$\\
若$f(x),g(x)$在$[a,b]$上Riemann可积且对任意$x\in [a,b]$,有$f(x)\leq g(x)$,则$\int_{a}^{b}f(x)\d x\leq \int_{a}^{b}g(x)\d x$
若$f(x),g(x)$在$[a,b]$上Riemann可积,对任意$\alpha,\beta\in\R$,则$\alpha f(x)+\beta g(x)$在$[a,b]$上Riemann可积\\
若$f(x),g(x)$在$[a,b]$上Riemann可积,则$f(x)g(x)$在$[a,b]$上Riemann可积
\end{dingli}

\begin{dingli}
对实数域$\R$\\
对函数$f,g:[a,b]\subset \R \to \R$,其中对任意$x\in[a,b]$,有$g(x)\geq 0$\\
若$f(x),g(x)$在$[a,b]$上Riemann可积且$f(x)$在$[a,b]$上连续\\
则存在$\xi\in[a,b]$,有$\int_{a}^{b}f(x)g(x)\d x=f(\xi)\int_{a}^{b}g(x)\d x$
\end{dingli}

\begin{dingli}
对实数域$\R$\\
对函数$f,g:[a,b]\subset \R \to \R$,其中对任意$x\in[a,b]$,有$f(x)\geq 0$\\
若$f(x)$在$[a,b]$上单调,则存在$\xi\in[a,b]$,有$\int_{a}^{b}f(x)g(x)\d x=f(a)\int_{a}^{\xi}g(x)\d x+f(b)\int_{\xi}^{b}g(x)\d x$\\
若$f(x)$在$[a,b]$上单调不增,则存在$\xi\in[a,b]$,有$\int_{a}^{b}f(x)g(x)\d x=f(a)\int_{a}^{\xi}g(x)\d x$\\
若$f(x)$在$[a,b]$上单调不减,则存在$\xi\in[a,b]$,有$\int_{a}^{b}f(x)g(x)\d x=f(b)\int_{\xi}^{b}g(x)\d x$
\end{dingli}

\begin{dingyi}
对实数域$\R$\\
对函数$f,g:D\subset \R \to \R$\\
若$f(x)$在$D$上可微且对任意$x\in D$,有$f'(x)=g(x)$\\
则称$g(x)$为$f(x)$在$D$上的原函数\\
则称$\int f(x)\d x=g(x)+C$为$f(x)$的不定积分,表示$f(x)$的原函数族,其中$C$为常数
\end{dingyi}

\begin{dingli}
微积分基本定理:\\
对实数域$\R$\\
对函数$f:[a,b]\subset \R \to \R$,其中$x_0\in[a,b]$\\
若$f(x)$在$[a,b]$上Riemann可积,则$F(x_0)=\int_{a}^{x_0}f(x)\d x$在$[a,b]$上连续\\
若$f(x)$在$x_0$处连续,则$F(x)$在$x_0$处可导或可微且$F'(x_0)=f(x_0)$\\
即$F(x)$为$f(x)$的原函数
\end{dingli}

\begin{dingli}
Newton-Leibniz:\\
对实数域$\R$\\
对函数$f:[a,b]\subset \R \to \R$\\
若$f(x)$在$[a,b]$上Riemann可积,则$\int_{a}^{b}f(x)\d x=(\int f(x)\d x)|_{x=a}^{b}$
\end{dingli}

\begin{zhu}
以下为常见函数的原函数:\\
\[\int x^{\alpha}\d x=\frac{x^{\alpha+1}}{\alpha+1}+C,\ x>0,\alpha\neq -1,\int \frac{1}{x}\d x=\ln|x|+C,\ x\neq0\] 
\[\int a^x\d x=\frac{x^{a}}{\ln a}+C, a<0,\int \e^x\d x=\e^x+C\]
\[\int \sin x\d x=-\cos x+C,\int\cos x\d x=\sin x+C\] 
\[\int \frac{1}{\cos^2 x}\d x=\tan x,\ x\neq \frac{\pi}{2}+n\pi,\int \frac{1}{\sin^2 x}\d x=-\cot x,\ x\neq n\pi\] 
\[\int \frac{1}{\sqrt{1-x^2}}\d x=\arcsin x+C,\ |x|<1,\int -\frac{1}{\sqrt{1-x^2}}\d x=\arccos x+C,\ |x|<1\]
\[\int \frac{1}{1+x^2}\d x=\arctan x+C\] 
\end{zhu}

\begin{dingli}
对实数域$\R$\\
对函数$f,g:D\subset \R \to \R$\\
若$F(x)$为$f(x)$在$D$上的原函数且$G(x)$为$g(x)$在$D$上的原函数\\
则$\int g(x)F(x)\d x=F(x)G(x)-\int f(x)G(x)\d x$
\end{dingli}

\begin{dingli}
对实数域$\R$\\
对函数$f:D\subset \R \to \R,g:D'\subset \R\to D$\\
若$g(x)$在$D'$上可微且$f(x)$在$D$上有原函数\\
则$\int f(x)\d x|_{x=g(y)}=\int f(g(y))g'(y)\d y$
\end{dingli}

\begin{dingli}
对实数域$\R$\\
对函数$f,g:[a,b]\subset \R \to \R$\\
若$F(x)$为$f(x)$在$[a,b]$上的原函数且$G(x)$为$g(x)$在$[a,b]$上的原函数\\
则$\int_{a}^{b} g(x)F(x)\d x=F(x)G(x)|_{x=a}^{b}-\int_{x=a}^{b} f(x)G(x)\d x$
\end{dingli}

\begin{dingli}
对实数域$\R$\\
对函数$f:[a,b]\subset \R \to \R,g:[c,d]\subset \R\to [a,b]$\\
若$g(x)$在$[c,d]$上可微且$f(x)$在$[a,b]$上有原函数且$g(c)=a,g(d)=b$\\
则$\int_{a}^{b} f(x)\d x=\int_{c}^{d}f(g(y))g'(y)\d y$
\end{dingli}

\begin{dingyi}
对实数域$\R$\\
对函数$f:[a,b)\subset \R \to \R$,其中$b=\{b_0\in\R|b_0>a\}\cup\{\infty\}$\\
若对任意$c\in(a,b)$,有$f(x)$在$[a,c]$上Riemann可积且$f(x)$在$b$任意领域无界或$|b-a|=\infty$\\
则称$\int_{a}^{b}f(x)\d x=\lim\limits_{c\to b^{-}}\int_{a}^{c}f(x)\d x$为$f(x)$在$[a,b)$上的反常积分\\
若极限$\lim\limits_{c\to b^{-}}\int_{a}^{c}f(x)\d x$存在\\
则称$\int_{a}^{b}f(x)\d x$为收敛的,否则称$\int_{a}^{b}f(x)\d x$为发散的\\
若$\int_{a}^{b}|f(x)|\d x$为收敛的\\
则称$\int_{a}^{b}f(x)\d x$为绝对收敛的\\
若$\int_{a}^{b}|f(x)|\d x$为发散的且$\int_{a}^{b}f(x)\d x$为收敛的\\
则称$\int_{a}^{b}f(x)\d x$为条件收敛的
\end{dingyi}

\begin{dingli}
Cauchy:\\
对实数域$\R$\\
对函数$f:[a,b)\subset \R \to \R$,其中$b=\{b_0\in\R|b_0>a\}\cup\{\infty\}$\\
若对任意$c\in(a,b)$,有$f(x)$在$[a,c]$上Riemann可积且$f(x)$在$b$任意领域无界或$|b-a|=\infty$\\
则$\int_{a}^{b}f(x)\d x$为收敛的当且仅当\\
对任意$\varepsilon>0$,存在$c\in (a,b)$,对任意$c'\in (c,b)$,有$|\int_{c}^{c'}f(x)\d x|<\varepsilon$
\end{dingli}

\begin{dingli}
Cauchy:\\
对实数域$\R$\\
对函数$f:[a,b)\subset(0,\infty)\subset \R \to \R$,其中$b=\{b_0\in\R|b_0>a\}\cup\{\infty\}$\\
若对任意$c\in(a,b)$,有$f(x)$在$[a,c]$上Riemann可积且$|b-a|=\infty$\\
若存在$c\in [a,b)$,对任意$x\in [e,b)$,有$0\leq f(x)\leq \frac{C}{x^p}$,其中$p>1,C>0$\\
则$\int_{a}^{b}f(x)\d x$为收敛的\\
若存在$c\in [a,b)$,对任意$x\in [e,b)$,有$f(x)\geq \frac{C}{x^p}$,其中$p\leq 1,C>0$\\
则$\int_{a}^{b}f(x)\d x$为发散的\\
若对任意$c\in(a,b)$,有$f(x)$在$[a,c]$上Riemann可积且$f(x)$在$b$任意领域无界\\
若存在$c\in [a,b)$,对任意$x\in [e,b)$,有$0\leq f(x)\leq \frac{C}{(b-x)^p}$,其中$p<1,C>0$\\
则$\int_{a}^{b}f(x)\d x$为收敛的\\
若存在$c\in [a,b)$,对任意$x\in [e,b)$,有$f(x)\geq \frac{C}{x^p}$,其中$p\geq 1,C>0$\\
则$\int_{a}^{b}f(x)\d x$为发散的
\end{dingli}

\begin{dingli}
Dirichlet:\\
对实数域$\R$\\
对函数$f,g:[a,b)\subset \R \to \R$,其中$b=\{b_0\in\R|b_0>a\}\cup\{\infty\}$\\
有对任意$c\in(a,b)$,有$f(x),g(x)$在$[a,c]$上Riemann可积且$f(x),g(x)$在$b$任意领域无界或$|b-a|=\infty$\\
若存在$M>0$,对任意$c\in(a,b)$,有$|\int_{a}^{c}f(x)\d x|<M$\\
且$g(x)$在$[a,b)$上单调,有$\lim\limits_{x\to b^{-}}g(x)=0$\\
则$\int_{a}^{b}f(x)g(x)\d x$为收敛的
\end{dingli}

\begin{dingli}
Abel:\\
对实数域$\R$\\
对函数$f,g:[a,b)\subset \R \to \R$,其中$b=\{b_0\in\R|b_0>a\}\cup\{\infty\}$\\
有对任意$c\in(a,b)$,有$f(x),g(x)$在$[a,c]$上Riemann可积且$f(x),g(x)$在$b$任意领域无界或$|b-a|=\infty$\\
若有$\int_{a}^{b}f(x)\d x$为收敛的且$g(x)$在$[a,b)$上单调有界\\
则$\int_{a}^{b}f(x)g(x)\d x$为收敛的
\end{dingli}

\begin{dingyi}
对$2$维实数域$\R^2$\\
对函数$f:[a,b]\times[c,d]\subset \R^2 \to \R$\\
若$f(x,y)$在$[a,b]\times[c,d]$上连续\\
则称$I(y)=\int_{a}^{b}f(x,y)\d x,J(x)=\int_{c}^{d}f(x,y)\d y$为含参变量积分
\end{dingyi}

\begin{dingli}
对$2$维实数域$\R^2$\\
对函数$f:[a,b]\times[c,d]\subset \R^2 \to \R$\\
若$f(x,y)$在$[a,b]\times[c,d]$上连续,则$I(y)$在$[c,d]$上连续且$J(x)$在$[a,b]$上连续\\
即对任意$y_0\in [c,d]$,有$\lim\limits_{y\to y_0}\int_{a}^{b}f(x,y)\d x=\int_{a}^{b}\lim\limits_{y\to y_0}f(x,y)\d x$\\
即对任意$x_0\in [a,b]$,有$\lim\limits_{x\to x_0}\int_{c}^{d}f(x,y)\d y=\int_{c}^{d}\lim\limits_{x\to x_0}f(x,y)\d y$
\end{dingli}

\begin{dingli}
对$2$维实数域$\R^2$\\
对函数$f:[a,b]\times[c,d]\subset \R^2 \to \R$\\
若$f(x,y)$在$[a,b]\times[c,d]$上可偏导且偏导数在$[a,b]\times[c,d]$上连续\\
则$\frac{\d}{\d y}\int_{a}^{b}f(x,y)\d x=I(y)'=\int_{a}^{b}\frac{\partial f}{\partial y}(x,y)\d x$\\
且$\frac{\d}{\d x}\int_{c}^{d}f(x,y)\d y=J(x)'=\int_{c}^{d}\frac{\partial f}{\partial x}(x,y)\d y$
\end{dingli}

\begin{dingli}
对$2$维实数域$\R^2$\\
对函数$f:[a,b]\times[c,d]\subset \R^2 \to \R$\\
若$f(x,y)$在$[a,b]\times[c,d]$上可偏导且偏导数在$[a,b]\times[c,d]$上连续\\
若有$g,h:[c,d]\to [a,b]$在$[c,d]$上可导或可微,记$F(y)=\int_{g(y)}^{h(y)}f(x,y)\d x$\\
则$F(y)$在$[c,d]$上可导或可微且$F'(y)=f(h(y),x)h'(y)-f(g(y),y)g'(y)+\int_{g(y)}^{h(y)}\frac{\partial f}{\partial y}(x,y)\d x$
\end{dingli}

\begin{dingyi}
对$2$维实数域$\R^2$\\
对函数$f:[a,b)\times[c,d]\subset \R^2 \to \R$,其中$b=\{b_0\in\R|b_0>a\}\cup\{\infty\}$\\
 若对任意$e\in(a,b)$,有$f(x,y)$在$[a,e]$关于$x$上Riemann可积且$f(x,y)$在$b$任意领域无界或$|b-a|=\infty$\\
则称$I(y)=\int_{a}^{b}f(x,y)\d x$为$f(x,y)$在$[a,b)$上的含参变量反常积分\\
若对任意$\varepsilon>0$,存在$e\in (a,b)$,对任意$y\in [c,d]$,有$|\int_{e}^{b}f(x,y)\d x|<\varepsilon$\\
则称$I(y)=\int_{a}^{b}f(x,y)\d x$在$[c,d]$上一致收敛
\end{dingyi}

\begin{dingli}
Cauchy:\\
对$2$维实数域$\R^2$\\
对函数$f:[a,b)\times[c,d]\subset \R^2 \to \R$,其中$b=\{b_0\in\R|b_0>a\}\cup\{\infty\}$\\
 若对任意$e\in(a,b)$,有$f(x,y)$在$[a,e]$关于$x$上Riemann可积且$f(x,y)$在$b$任意领域无界或$|b-a|=\infty$\\
则$I(y)=\int_{a}^{b}f(x,y)\d x$在$[c,d]$上一致收敛当且仅当\\
对任意$\varepsilon>0$,存在$e\in (a,b)$,对任意$e'\in(e,b),y\in [c,d]$,有$|\int_{e}^{e'}f(x,y)|<\varepsilon$
\end{dingli}

\begin{dingli}
Weierstrass:\\
对$2$维实数域$\R^2$\\
对函数$f:[a,b)\times[c,d]\subset \R^2 \to \R$,其中$b=\{b_0\in\R|b_0>a\}\cup\{\infty\}$\\
若对任意$e\in(a,b)$,有$f(x,y)$在$[a,e]$关于$x$上Riemann可积且$f(x,y)$在$b$任意领域无界或$|b-a|=\infty$\\
若存在$F:[a,b)\subset \R\to\R$,对任意$(x,y)\in[a,b)\times[c,d]$,有$|f(x,y)|\leq F(x)$且$\int_{a}^{b}F(x)\d x$收敛\\
则$I(y)=\int_{a}^{b}f(x,y)\d x$在$[c,d]$上一致收敛
\end{dingli}

\begin{dingli}
Dirichlet:\\
对$2$维实数域$\R^2$\\
对函数$f:[a,b)\times[c,d]\subset \R^2 \to \R$,其中$b=\{b_0\in\R|b_0>a\}\cup\{\infty\}$\\
若对任意$e\in(a,b)$,有$f(x,y),g(x,y)$在$[a,e]$关于$x$上Riemann可积且$f(x,y),g(x,y)$在$b$任意领域无界或$|b-a|=\infty$\\
若存在$M>0$,对任意$e\in (a,b),y\in [c,d]$,有$|\int_{a}^{e}f(x,y)\d x|\leq M$\\
若对任意$y\in [c,d]$,有$g(x,y)$关于$x$单调\\
若对任意$\varepsilon>0$,存在$e\in (a,b)$,对任意$x\in(e,b),y\in [c,d]$,有$|f(x,y)|<\varepsilon$\\
则$I(y)=\int_{a}^{b}f(x,y)g(x,y)\d x$在$[c,d]$上一致收敛
\end{dingli}

\begin{dingli}
Abel:\\
对$2$维实数域$\R^2$\\
对函数$f:[a,b)\times[c,d]\subset \R^2 \to \R$,其中$b=\{b_0\in\R|b_0>a\}\cup\{\infty\}$\\
若对任意$e\in(a,b)$,有$f(x,y),g(x,y)$在$[a,e]$关于$x$上Riemann可积且$f(x,y),g(x,y)$在$b$任意领域无界或$|b-a|=\infty$\\
若$\int_{a}^{b}f(x,y)\d x$在$[c,d]$上一致收敛\\
若对任意$y\in [c,d]$,有$g(x,y)$关于$x$单调\\
若存在$M>0$,对任意$(x,y)\in[a,b)\times[c,d]$,有$|f(x,y)|<M$\\
则$I(y)=\int_{a}^{b}f(x,y)g(x,y)\d x$在$[c,d]$上一致收敛
\end{dingli}

\begin{zhu}
以下为Euler积分:\\
Gamma:$\Gamma(s)=\int_{0}^{\infty}x^{s-1}\e^{-x}\d x,\ s\in (0,\infty)$\\
有$\Gamma(s)$在$(0,\infty)$上可微或可导\\
有$\Gamma(s+1)=s\Gamma(s)$且$\Gamma(s)\Gamma(s+\frac{1}{2})=2^{1-2s}\sqrt{\pi}\Gamma(2s)$\\
有$\Gamma(s)\Gamma(1-s)=\frac{\pi}{\sin\pi s}$且$\Gamma(\frac{1}{2})=\sqrt{\pi}$\\
Beta:$B(p,q)=\int_{0}^{1}x^{p-1}(1-x)^{q-1}\d x,\ (p,q)\in(0,\infty)\times(0,\infty)$\\
有$B(p,q)$在$(0,\infty)\times(0,\infty)$上连续\\
有$B(p,q)=B(q,p)$且$B(p,q+1)=\frac{q}{p+q}B(p,q)$\\
有$B(\frac{1}{2},\frac{1}{2})=\pi$\\
有$B(p,q)=\frac{\Gamma(p)\Gamma(p)}{\Gamma(p+q)}$
\end{zhu}

\chapter{多元}

\begin{dingyi}
对实数域$\R$\\
记$\R^n=\R\times \R\times \cdots \times \R=\{(x_1,\cdots,x_n)|x_i\in\R,i=1,\cdots,n\}$,其中$n\in\N^*$\\
则称$\R^n$为$n$维实数域\\
对任意$x,y\in\R^n$,记$|x-y|=\sqrt{(x_1-y_1)^2+\cdots+(x_n-y_n)^2}$\\
则称$|x-y|$为$\R^n$上的距离\\
对子集$D\subset \R^n$,若存在$M>0$\\
对任意$x\in D$,有$|x|<M$,则称$D$为有界的
\end{dingyi}

\begin{dingyi}
对$n$维实数域$\R^n$\\
对数列$\{x^{(k)}\}\subset \R^n$,若存在$x\in \R^n$\\
对任意$\varepsilon>0$,存在$N\in \N$,对任意$k\geq N$,有$|x^{(k)}-x|<\varepsilon$\\
则称$\{x^{(k)}\}$为收敛的,记为$\lim\limits_{k\to\infty}x^{(k)}=x$,否则称$\{x^{(k)}\}$为发散的
\end{dingyi}

\begin{dingli}
对$n$维实数域$\R^n$\\
则$\{x^{(k)}\}$为收敛的当且仅当对任意$i=1,\cdots,n$,有$\{x_i^{(k)}\}$为收敛的
\end{dingli}

\begin{dingyi}
对$n$维实数域$\R^n$\\
对任意$x_0\in \R^n,\varepsilon$,记$O(x_0,\varepsilon)=\{x\in\R^n\Big||x-x_0|<\varepsilon\}$\\
则称$O(x_0,\varepsilon)$为$x_0$的邻域\\
对子集$D\subset \R^n$,若对$x\in \R^n$,存在$\varepsilon>0$,有$O(x,\varepsilon)\subset D$,则称$x$为$D$的内点\\
对子集$D\subset \R^n$,若对$x\in \R^n$,存在$\varepsilon>0$,有$O(x,\varepsilon)\subset D^{c}$,则称$x$为$D$的外点\\
对子集$D\subset \R^n$,若对$x\in \R^n$,对任意$\varepsilon>0$,有$O(x,\varepsilon)\cap D\neq\varnothing$且$O(x,\varepsilon)\cap D^{c}\neq \varnothing$\\
则称$x$为$D$的界点\\
对子集$D\subset \R^n$,若对$x\in \R^n$,对任意$\varepsilon>0$,有$O(x,\varepsilon)\cap D$包含$D$中$\infty$个点\\
则称$x$为$D$的聚点\\
对子集$D\subset \R^n$,记$\overline{D}=\{x\in\R^n|x\text{为}D\text{的聚点}\}\cup D$,则称$\overline{D}$为$D$的闭包\\
对子集$D\subset \R^n$,若$\{x\in\R^n|x\text{为}D\text{的内点}\}=D$,则称$D$为开集\\
对子集$D\subset \R^n$,若$\{x\in\R^n|x\text{为}D\text{的聚点}\}\subset D$,则称$D$为闭集
\end{dingyi}

\begin{dingli}
对$n$维实数域$\R^n$\\
对子集$D\subset \R^n$,其中$x\in \R^n$\\
则$O(x,\varepsilon)$为开集,则$\overline{D}$为闭集\\
则$x$为$D$的聚点当且仅当存在$\{x^{(k)}\}\subset D$,有$\lim\limits_{k\to\infty}x^{(k)}=x$
\end{dingli}

\begin{dingli}
Cantor:\\
对$n$维实数域$\R^n$\\
若有闭集列$D_1\supset D_2\supset\cdots\supset D_k\supset\cdots$,记$d(D_k)=\sup\{|x-y|\Big|x,y\in D_k\}$\\
若有$\lim\limits_{k\to\infty}d(D_k)=0$,则存在唯一$x\in \R^n$,有$x\in \bigcap\limits_{k=1}^{\infty}D_k$
\end{dingli}

\begin{dingli}
Heine-Borel:\\
对$n$维实数域$\R^n$\\
对任意有界闭集$D$和开集族$\{D_{i}|i\in I\}$\\
若$D\subset \bigcup\limits_{i\in I}D_{i}$,则存在$J\subset I,|J|<\infty$,有$D\subset \bigcup\limits_{j\in J}D_{j}$\\
即$D$为紧集
\end{dingli}

\begin{dingli}
Bolzano-Weierstrass:\\
对$n$维实数域$\R^n$\\
对子集$D\subset \R,|D|=\infty$\\
若$D$为有界的,则存在$x\in \R^n$为$D$的聚点\\
对数列$\{x^{(k)}\}$,若$\{x^{(k)}\}$为有界的\\
则存在子列$\{x^{(k_i)}\}\subset\{x^{(k)}\},x\in\R^n$,有$\lim\limits_{i\to\infty}x^{(k_{i})}=x$
\end{dingli}

\begin{dingli}
Cauchy:\\
对$n$维实数域$\R^n$\\
对数列$\{x^{(k)}\}$,若$\{x^{(k)}\}$为基本列\\
即对任意$\varepsilon>0$,存在$N\in \N$,对任意$k,l\geq N$,有$|x^{(k)}-x^{(l)}|<\varepsilon$\\
则存在$x\in \R^n$,有$\lim\limits_{k\to\infty}x^{(k)}=x$
\end{dingli}

\begin{dingyi}
对$n$维实数域$\R^n$\\
则称映射$f:D\subset \R^n\to \R$为函数,则称$D$为$f$的定义域,则称$f(D)$为值域\\
若$f(D)$为有界的,则称$f(x)$在$D$上有界\\
若存在$\alpha\in\R,a\in \R^n$,对任意$\varepsilon>0$,存在$\delta>0$,对任意$0<|x-a|<\delta$,有$|f(x)-\alpha|<\varepsilon$\\
则称$x$趋于$a$时$f(x)$有极限$\alpha$,记为$\lim\limits_{x\to a}f(x)=\alpha$\\
若对$x=(x_1,\cdots,x_n),a=(a_1,\cdots,a_n)$,有$k_1,\cdots,k_n$为$1,\cdots,n$的排列\\
则称$\lim\limits_{x_{k_1}\to a_{k_1}}\cdots\lim\limits_{x_{k_n}\to a_{k_n}}f(x_1,\cdots,x_n)$为累次极限
\end{dingyi}

\begin{dingli}
对$2$维实数域$\R^2$\\
对函数$f:D\subset\R^2\to\R$,其中$x^{(0)}=(x_1^{(0)},x_2^{(0)})\in\R^2$\\
若极限$\lim\limits_{x\to x^{(0)}}f(x)$存在且对任意$x_1\neq x_1^{(0)}$,极限$\lim\limits_{x_2\to x_2^{(0)}}f(x_1,x_2)=g(x_1)$存在\\
则累次极限$\lim\limits_{x_1\to x_1^{(0)}}\lim\limits_{x_2\to x_2^{(0)}}f(x_1,x_2)$存在且$\lim\limits_{x_1\to x_1^{(0)}}\lim\limits_{x_2\to x_2^{(0)}}f(x_1,x_2)=\lim\limits_{x\to x^{(0)}}f(x)$\\
若累次极限$\lim\limits_{x_1\to x_1^{(0)}}\lim\limits_{x_2\to x_2^{(0)}}f(x_1,x_2)$和$\lim\limits_{x_2\to x_2^{(0)}}\lim\limits_{x_1\to x_1^{(0)}}f(x_1,x_2)$和极限$\lim\limits_{x\to x^{(0)}}f(x)$均存在\\
则$\lim\limits_{x_1\to x_1^{(0)}}\lim\limits_{x_2\to x_2^{(0)}}f(x_1,x_2)=\lim\limits_{x_2\to x_2^{(0)}}\lim\limits_{x_1\to x_1^{(0)}}f(x_1,x_2)=\lim\limits_{x\to x^{(0)}}f(x)$
\end{dingli}

\begin{dingyi}
对$n$维实数域$\R^n$\\
对函数$f:D\subset \R^n \to \R$,其中$x^{(0)}\in D$\\
若对任意$\varepsilon>0$,存在$\delta>0$,对任意$x\in D,|x-x^{(0)}|<\delta$,有$|f(x)-f(x^{(0)})|<\varepsilon$\\
则称$f(x)$在$x_0$处为连续的\\
若对任意$\varepsilon>0$,存在$\delta>0$,对任意$x^{(1)},x^{(2)}\in D,|x^{(1)}-x^{(2)}|<\delta$,有$|f(x^{(1)})-f(x^{(2)})|<\varepsilon$\\
则称$f(x)$在$D$上一致连续
\end{dingyi}

\begin{dingli}
对$n$维实数域$\R^n$\\
对函数$f:D\subset \R^n \to \R$\\
若$D$为紧集且$f(x)$在$D$上连续,则$f(x)$在$D$上有界\\
且存在$x^{(1)},x^{(2)}\in D$,对任意$x\in D$,有$f(x^{(1)})\leq f(x)\leq f(x^{(2)})$\\
则$f(x)$在$D$上一致连续
\end{dingli}

\begin{dingyi}
对$n$维实数域$\R^n$\\
对函数$f:D\subset \R^n \to \R$,其中$x^{(0)}=(x_1^{(0)},\cdots,x_n^{(0)})\in D$\\
若极限$\lim\limits_{h\to 0}\frac{f(x_1^{(0)},\cdots,x_{i-1}^{(0)},x_i^{(0)}+h,x_{i+1}^{(0)},\cdots,x_n^{(0)})-f(x^{(0)})}{h}$存在\\
则称$f(x)$在$x^{(0)}$处关于$x_i$可偏导,记为$\frac{\partial f}{\partial x_i}|_{x=x^{(0)}}$\\
若对任意$x^{(0)}\in D$,有$f(x)$在$x^{(0)}$处关于$x_i$可偏导\\
则称$f(x)$在$D$上关于$x_i$可偏导\\
则称$\frac{\partial f}{\partial x_i}:D\to \R,x^{(0)}\to \frac{\partial f}{\partial x_i}|_{x=x^{(0)}}$为$f(x)$在$D$上关于$x_i$的偏导函数\\
若对任意$i=1,\cdots,n$,有$f(x)$在$D$上关于$x_i$可偏导\\
则称$f(x)$在$D$上可偏导
\end{dingyi}

\begin{dingyi}
对$n$维实数域$\R^n$\\
对函数$f:D\subset \R^n \to \R$,其中$x^{(0)}=(x_1^{(0)},\cdots,x_n^{(0)})\in D$\\
若有方向向量$v=(v_1,\cdots,v_n)\in\R^n$,其中$v_1^2+\cdots,+v_n^2=1$\\
则称$\lim\limits_{h\to 0}\frac{f(x_1^{(0)}+hv_1,\cdots,x_n^{(0)}+hv_n)-f(x^{(0)})}{h}$为$f(x)$在$x^{(0)}$处关于$v$的方向导数,记为$\frac{\partial f}{\partial v}|_{x=x^{(0)}}$
\end{dingyi}

\begin{dingyi}
对$n$维实数域$\R^n$\\
对函数$f:D\subset \R^n \to \R$,其中$x^{(0)}=(x_1^{(0)},\cdots,x_n^{(0)})\in D$\\
记$\Delta x_i$,记$\Delta f=f(x_1^{(0)}+\Delta x_1,\cdots,x_n^{(0)}+\Delta x_n)-f(x^{(0)})$\\
若存在与$\Delta x_i$无关的常数$A_1,\cdots,A_n$,有$\Delta f=A_1\Delta x_1+\cdots+A_n\Delta x_n+o(\sqrt{\Delta x_1^2+\cdots+\Delta x_n^2})$\\
则称$f(x)$在$x^{(0)}$处可微\\
则称$\d f|_{x=x^{(0)}}=A_1\d x_1+\cdots+A_n\d x_n$为$f(x)$在$x^{(0)}$处的全微分\\
若对任意$x^{(0)}\in D$,有$f(x)$在$x^{(0)}$处可微\\
则称$f(x)$在$D$上可微\\
则称$\d f:D\to \R\d x,x^{(0)}\to A_1\d x_1+\cdots+A_n\d x_n$为$f(x)$在$D$上的全微分形式
\end{dingyi}

\begin{dingli}
对$n$维实数域$\R^n$\\
对函数$f:D\subset \R^n \to \R$,其中$x^{(0)}=(x_1^{(0)},\cdots,x_n^{(0)})\in D$\\
若$f(x)$在$x^{(0)}$处可微,则$f(x)$在$x^{(0)}$处连续且可偏导\\
有$\d f|_{x=x^{(0)}}=\frac{\partial f}{\partial x_1}|_{x=x^{(0)}}\d x_1+\cdots+\frac{\partial f}{\partial x_n}|_{x=x^{(0)}}\d x_n$\\
且对任意方向向量$v=(v_1,\cdots,v_n)\in\R^n$,其中$v_1^2+\cdots,+v_n^2=1$
有方向导数存在且$\lim\limits_{h\to 0}\frac{f(x_1^{(0)}+hv_1,\cdots,x_n^{(0)}+hv_n)-f(x^{(0)})}{h}=\frac{\partial f}{\partial x_1}|_{x=x^{(0)}}h_1+\cdots+\frac{\partial f}{\partial x_n}|_{x=x^{(0)}}h_n$\\
若$f(x)$在$D$上可微,则$f(x)$在$D$上连续且可偏导\\
有$\d f=\frac{\partial f}{\partial x_1}\d x_1+\cdots+\frac{\partial f}{\partial x_n}\d x_n$\\
若存在$O(x^{(0)},\varepsilon)$,有$f(x)$在$O(x^{(0)},\varepsilon)$上可偏导且偏导数在$x^{(0)}$处连续\\
则$f(x)$在$x^{(0)}$处可微
\end{dingli}

\begin{dingli}
对$2$维实数域$\R^2$\\
对函数$f:D\subset\R^2\to\R$,其中$x_1=x_1(u,v),x_2=x_2(u,v):D'\to D$\\
若对$x^{(0)}=(x_1^{(0)},x_2^{(0)})\in D,(u_0,v_0)\in D'$,有$x_1^{(0)}=x_1(u_0,v_0),x_2^{(0)}=x_2(u_0,v_0)$\\
且$x_1(u,v),x_2(u,v)$在$(u_0.v_0)$处可偏导且$f(x)$在$x^{(0)}$处可微\\
则有$\frac{\partial f}{\partial u}|_{(u_0,v_0)}=\frac{\partial f}{\partial x_1}|_{x=x^{(0)}}\frac{\partial x_1}{\partial u}|_{(u_0.v_0)}+\frac{\partial f}{\partial x_2}|_{x=x^{(0)}}\frac{\partial x_2}{\partial u}|_{(u_0.v_0)}$\\
则有$\frac{\partial f}{\partial v}|_{(u_0,v_0)}=\frac{\partial f}{\partial x_1}|_{x=x^{(0)}}\frac{\partial x_1}{\partial v}|_{(u_0.v_0)}+\frac{\partial f}{\partial x_2}|_{x=x^{(0)}}\frac{\partial x_2}{\partial v}|_{(u_0.v_0)}$
\end{dingli}

\begin{dingli}
对$n$维实数域$\R^n$\\
对函数$f:D\subset \R^n \to \R$\\
若$D$为凸区域且$f(x)$在$D$上可微,则对任意$x^{(1)},x^{(2)}\in D$,存在$\theta\in (0,1)$,有$x^{(0)}=x^{(1)}+\theta(x^{(2)}-x^{(1)})$\\
有$f(x_1^{(1)},x_2^{(1)})-f(x_1^{(2)},x_2^{(2)})=\frac{\partial f}{\partial x_1}|_{x=x^{(0)}}(x_1^{(1)}-x_1^{(2)})+\frac{\partial f}{\partial x_2}|_{x=x^{(0)}}(x_2^{(1)}-x_2^{(2)})$
\end{dingli}

\begin{dingyi}
对$n$维实数域$\R^n$\\
对函数$f:D\subset \R^n \to \R$\\
若$f(x)$在$D$上可偏导,则可递归定义偏导函数的偏导函数,记$i_1,\cdots,i_m\in\{1,\cdots,n\}$\\
则称$\frac{\partial^m f}{\partial x_{i_m}\cdots\partial x_{i_1}}=\frac{\partial}{\partial x_{i_m}}(\cdots(\frac{\partial f}{\partial x_{i_1}}))$为$f(x)$的$m$阶偏导数\\
若对任意$i_1,\cdots,i_m\in\{1,\cdots,n\}$,偏导函数$\frac{\partial^m f}{\partial x_{i_m}\cdots\partial x_{i_1}}=\frac{\partial}{\partial x_{i_m}}(\cdots(\frac{\partial f}{\partial x_{i_1}}))$存在\\
则称$f(x)$在$D$上可$m$阶偏导
\end{dingyi}

\begin{dingli}
对函数$f:D\subset \R^n \to \R$\\
若$f(x)$在$D$上可$2$阶偏导,且$\frac{\partial^2 f}{\partial x_i\partial x_j}$和$\frac{\partial^2 f}{\partial x_j\partial x_i}$均在$D$上连续\\
则有$\frac{\partial^2 f}{\partial x_i\partial x_j}=\frac{\partial^2 f}{\partial x_j\partial x_i}$
\end{dingli}

\begin{dingyi}
对$n$维实数域$\R^n$\\
对函数$f:D\subset \R^n \to \R$\\
若$f(x)$在$D$上可微,则可递归定义微分形式的微分\\
则称$\d^m f=\d(\d^{m-1}f)$为$f(x)$的$m$阶微分形式
\end{dingyi}

\begin{zhu}
由数学归纳法和$\d^{m} x_i=0,m\geq 2$可知\\
则有$\d^m f=(\frac{\partial }{\partial x_1}\d x_1+\cdots+\frac{\partial }{\partial x_n}\d x_n)^{m}f$
\end{zhu}

\begin{dingli}
对$n$维实数域$\R^n$\\
对函数$f:D\subset \R^n \to \R$,其中$x^{(0)}=(x_1^{(0)},\cdots,x_n^{(0)})\in D$\\
若$f(x)$在$O(x^{(0)},\varepsilon)$上可$m+1$阶偏导且偏导数连续\\
则对任意$x^{(0)}+\Delta x\in O(x^{(0)},\varepsilon)$\\
有$f(x^{(0)}+\Delta x)=f(x^{(0)})+(\sum\limits_{i=1}^{n}\frac{\partial }{\partial x_i}\Delta x_i)f|_{x^{(0)}}+\frac{1}{2!}(\sum\limits_{i=1}^{n}\frac{\partial }{\partial x_i}\Delta x_i)^2f|_{x^{(0)}}+\cdots+\frac{1}{m!}(\sum\limits_{i=1}^{n}\frac{\partial }{\partial x_i}\Delta x_i)^mf|_{x^{(0)}}+T_m$\\
其中$T_n=\frac{1}{m+1}(\sum\limits_{i=1}^{n}\frac{\partial }{\partial x_i}\Delta x_i)^{m+1}f|_{x^{(0)}+\theta\Delta x}$且$\theta\in(0,1)$
\end{dingli}

\begin{dingyi}
对$n$维实数域$\R^n$\\
对函数$f:D\subset \R^n \to \R$\\
若有分划$D=\bigcup\limits_{i=1}^{n}D_i$且对$i\neq j$,有$\Delta (D_i\cap D_j)=0$\\
其中$\Delta D$为测度意义下的积,$(\ n=2$为面积;$n=3$为体积$)$\\
若有$\xi^{(i)}\in D_i,i=1,2,\cdots,n$,记$\lambda=\max\limits_{i}\{d(D_i)\}$,其中$d(D_i)=\sup\{|x-y|\Big|x,y\in D_i\}$\\
则称$\sum\limits_{i=1}^{n}f(\xi^{(i)})\Delta D_i$为$f(x)$关于此分划的Riemann和\\
若$\lim\limits_{\lambda\to 0}\sum\limits_{i=1}^{n}f(\xi^{(i)})\Delta D_i$存在且无关$\xi^{(i)}$的选取\\
则称$\int_{D}f(x)\d \Delta=\lim\limits_{\lambda\to 0}\sum\limits_{i=1}^{n}f(\xi^{(i)})\Delta D_i$为$f(x)$在$D$上的$m$重积分\\
则称$f(x)$在$D$上可积
\end{dingyi}

\begin{dingli}
对$n$维实数域$\R^n$\\
对函数$f:D\subset \R^n \to \R$\\
若$D=[a_1,b_1]\times[a_n,b_n]$且$f(x)$在$D$上连续\\
则$\int_{D}f(x)\d \Delta=\int_{a_1}^{b_1}\cdots\int_{a_n}^{b_n}f(x_1,\cdots,x_n)\d x_1\cdots\d x_n$
\end{dingli}

\begin{dingli}
对$n$维实数域$\R^n$\\
对函数$f:D\subset \R^n \to \R$\\
若$D=[a_1,b_1]\times[a_n,b_n]$且$f(x)$在$D$上可积,记$D'=[a_2,b_2]\times[a_n,b_n]$\\
若对任意$x_1\in [a_1,b_1]$,有$g(x_1)=\int_{D'}f(x')\d \Delta'$存在\\
则$g(x_1)$在$[a_1,b_1]$上Riemann可积且$\int_{a}^{b}g(x_1)\d x_1=\int_{D}f(x)\d \Delta$
\end{dingli}

\begin{dingli}
对$n$维实数域$\R^n$\\
对函数$f:D\subset \R^n \to \R$\\
若有一一映射$T:D'\subset \R^n\to D,y\to x$,则$T(y_1,\cdot,y_n)=(x_1(y_1,\cdot,y_n),\cdots,x_n(y_1,\cdot,y_n))$\\
若$x_i(y)$在$D'$上可偏导且偏导数在$D'$上连续且$|J|=|(\frac{\partial x_i}{\partial y_j})|\neq 0$\\
若$f(x)$在$D$上可积,则$\int_{D}f(x)\d\Delta =\int_{D'}f(T(y))|J|\d \Delta'$
\end{dingli}

\begin{dingyi}
对$n$维实数域$\R^n$\\
对函数$f:D\subset \R^n \to \R$\\
在$D$内有一可求长曲线$L$,有划分$P_0\cdots P_n$\\
若有$x^{(i)}\in P_{i-1}P_{i}$且$\Delta s_i=P_{i-1}P_{i}$,记$\lambda=\sup\limits_{i}\{\Delta s_{i}\}$\\
若$\lim\limits_{\lambda\to 0}\sum\limits_{i=1}^{n}f(x^{(i)})\Delta s_i$存在且无关$x^{(i)}$的选取\\
则称$\int_{L}f(x)\d s=\lim\limits_{\lambda\to 0}\sum\limits_{i=1}^{n}f(x^{(i)})\Delta s_i$为第一类曲线积分\\
在$D$内有一可求长曲线$L$且$f(x)=\sum\limits_{i=1}^{n}f_i(x)\vec{x_i}$\\
若在$L$上每点有方向向量$v_{L}=(v_1,\cdots,v_n)\in\R^n$,其中$v_1^2+\cdots,+v_n^2=1$\\
则称$\int_{L}f(x)\cdot v_{L}\d s=\int_{L}\sum\limits_{i=1}^{n}f_{i}(x)v_i\d s$为第二类曲线积分
\end{dingyi}

\begin{dingli}
对$n$维实数域$\R^n$\\
对函数$f:D\subset \R^n \to \R$\\
在$D$内有一可求长光滑曲线$L$且有$x(t):[a,b]\to L$\\
若$f(x)$在$L$上连续,则$\int_{L}f(x)\d s=\int_{a}^{b}f(x(t))\sqrt{\sum\limits_{i=1}^{n}x_i'^2(t)}\d t$
\end{dingli}

\begin{dingyi}
对$n$维实数域$\R^n$\\
对函数$f:D\subset \R^n \to \R$\\
在$D$内有一可求积曲面$\sum$,有划分$\sum=\bigcup\limits_{i=1}^{n}\sum_{i}$且对$i\neq j$,有$\Delta (\sum_i\cap \sum_j)=0$\\
若有$x^{(i)}\in \sum_{i}$且$\Delta S_i=\Delta\sum_i$,记$\lambda=\sup\limits_{i}\{d(\sum_i)\}$\\
若$\lim\limits_{\lambda\to 0}\sum\limits_{i=1}^{n}f(x^{(i)})\Delta S_i$存在且无关$x^{(i)}$的选取\\
则称$\iint_{\sum}f(x)\d S=\lim\limits_{\lambda\to 0}\sum\limits_{i=1}^{n}f(x^{(i)})\Delta S_i$为第一类曲线积分\\
在$D$内有一可求积定向曲面$\sum$且$f(x)=\sum\limits_{i=1}^{n}f_i(x)\vec{x_i}$\\
若在$\sum$上每点有法向量$v_{\sum}=(v_1,\cdots,v_n)\in\R^n$,其中$v_1^2+\cdots,+v_n^2=1$\\
则称$\iint_{\sum}f(x)\cdot v_{\sum}\d S=\iint_{\sum}\sum\limits_{i=1}^{n}f_{i}(x)v_i\d S$为第二类曲线积分
\end{dingyi}

\begin{dingli}
Green:\\
对$2$维实数域$\R^2$\\
对函数$f,g:D\subset \R^2 \to \R$,其中$D$为单连通区域且$\partial D$为分段光滑简单闭曲线\\
若$f(x),g(x)$在$D$上可偏导且偏导数在$D$上连续,取$\partial D$的正向\\
则有$\int_{\partial D}f(x)\d x_1+g(x)\d x_2=\int_{D}(\frac{\partial g }{\partial x_1}-\frac{\partial f }{\partial x_2})\d x_1\d x_2$
\end{dingli}

\begin{dingli}
Gauss:\\
对$3$维实数域$\R^3$\\
对函数$f,g,h:D\subset \R^3 \to \R$,其中$D$为单连通区域且$\partial D$为分段光滑简单闭曲面\\
若$f(x),g(x),h(x)$在$D$上可偏导且偏导数在$D$上连续,取$\partial D$的正向\\
则有$\iint_{\partial D}f(x)\d x_2\d x_3+g(x)\d x_3\d x_1+h(x)\d x_1\d x_2=\int_{D}(\frac{\partial f }{\partial x_1}+\frac{\partial g }{\partial x_2}+\frac{\partial h }{\partial x_3})\d x_1\d x_2\d x_3$
\end{dingli}

\begin{dingli}
Stokes:\\
对$3$维实数域$\R^3$\\
对函数$f,g,h:D\subset \R^3 \to \R$,其中$\sum\subset D$为光滑曲面且$\partial\sum$为分段光滑简单闭曲线\\
若$f(x),g(x),h(x)$在$\sum\cup\partial\sum$上可偏导且偏导数在$\sum\cup\partial\sum$上连续,取$\partial \sum$的正向\\
则有$\int_{\partial \sum}f(x)\d x_1+g(x)\d x_2+h(x)\d x_3$\\
\hspace*{2em}$=\iint_{\sum}(\frac{\partial h }{\partial x_2}-\frac{\partial g }{\partial x_3})\d x_2\d x_3+(\frac{\partial f}{\partial x_3}-\frac{\partial h }{\partial x_1})\d x_3\d x_1+(\frac{\partial g }{\partial x_1}-\frac{\partial f}{\partial x_2})\d x_1\d x_2$
\end{dingli}

\endinput
%\end{document}